% Ngân hàng đề Toán 12 (KNTT) -- Chương II, Bài 8, Dạng 11
% Dạng: Ứng dụng thực tế: định vị, bay, cáp treo
%% Phương án ĐẦU TIÊN trong \cauTN luôn là đáp án đúng (gói trondethi tự đảo).

\cauTN[id=C2-B8-D11-NB-001,muc=NB,loigiai={Trạm trung chuyển đặt tại trung điểm $I$ của $AB$: $I=(2;\dfrac{1}{2};\dfrac{1}{2})$.}]{Trong không gian $Oxyz$ (đơn vị: km), hai trạm phát sóng đặt tại $A(-1;-2;-3)$ và $B(5;3;4)$. Một trạm trung chuyển được đặt tại trung điểm của đoạn thẳng nối $A$ và $B$. Tọa độ của trạm trung chuyển là}
  {$(2;\dfrac{1}{2};\dfrac{1}{2})$.}
  {$(2;\dfrac{3}{2};\dfrac{1}{2})$.}{$(2;\dfrac{5}{2};\dfrac{1}{2})$.}{$(-2;-\dfrac{1}{2};-\dfrac{1}{2})$.}

\cauTN[id=C2-B8-D11-NB-002,muc=NB,loigiai={Trạm trung chuyển đặt tại trung điểm $I$ của $AB$: $I=(-\dfrac{3}{2};4;-1)$.}]{Trong không gian $Oxyz$ (đơn vị: km), hai trạm phát sóng đặt tại $A(-2;4;1)$ và $B(-1;4;-3)$. Một trạm trung chuyển được đặt tại trung điểm của đoạn thẳng nối $A$ và $B$. Tọa độ của trạm trung chuyển là}
  {$(-\dfrac{3}{2};4;-1)$.}
  {$(-\dfrac{3}{2};4;-2)$.}{$(1;0;-4)$.}{$(\dfrac{3}{2};-4;1)$.}

\cauTN[id=C2-B8-D11-NB-003,muc=NB,loigiai={Trạm trung chuyển đặt tại trung điểm $I$ của $AB$: $I=(-2;1;-\dfrac{1}{2})$.}]{Trong không gian $Oxyz$ (đơn vị: km), hai trạm phát sóng đặt tại $A(-5;-3;0)$ và $B(1;5;-1)$. Một trạm trung chuyển được đặt tại trung điểm của đoạn thẳng nối $A$ và $B$. Tọa độ của trạm trung chuyển là}
  {$(-2;1;-\dfrac{1}{2})$.}
  {$(-\dfrac{1}{2};1;-2)$.}{$(6;8;-1)$.}{$(-2;-\dfrac{1}{2};1)$.}

\cauTN[id=C2-B8-D11-NB-004,muc=NB,loigiai={Trạm trung chuyển đặt tại trung điểm $I$ của $AB$: $I=(3;\dfrac{3}{2};\dfrac{3}{2})$.}]{Trong không gian $Oxyz$ (đơn vị: km), hai trạm phát sóng đặt tại $A(4;3;4)$ và $B(2;0;-1)$. Một trạm trung chuyển được đặt tại trung điểm của đoạn thẳng nối $A$ và $B$. Tọa độ của trạm trung chuyển là}
  {$(3;\dfrac{3}{2};\dfrac{3}{2})$.}
  {$(\dfrac{3}{2};\dfrac{3}{2};3)$.}{$(3;\dfrac{1}{2};\dfrac{3}{2})$.}{$(3;\dfrac{3}{2};-\dfrac{3}{2})$.}

\cauTN[id=C2-B8-D11-NB-005,muc=NB,loigiai={Trạm trung chuyển đặt tại trung điểm $I$ của $AB$: $I=(-\dfrac{3}{2};-2;\dfrac{7}{2})$.}]{Trong không gian $Oxyz$ (đơn vị: km), hai trạm phát sóng đặt tại $A(-2;-2;4)$ và $B(-1;-2;3)$. Một trạm trung chuyển được đặt tại trung điểm của đoạn thẳng nối $A$ và $B$. Tọa độ của trạm trung chuyển là}
  {$(-\dfrac{3}{2};-2;\dfrac{7}{2})$.}
  {$(\dfrac{1}{2};-2;\dfrac{7}{2})$.}{$(-\dfrac{3}{2};-3;\dfrac{7}{2})$.}{$(\dfrac{3}{2};-2;\dfrac{7}{2})$.}

\cauTN[id=C2-B8-D11-NB-006,muc=NB,loigiai={Trạm trung chuyển đặt tại trung điểm $I$ của $AB$: $I=(0;\dfrac{3}{2};\dfrac{1}{2})$.}]{Trong không gian $Oxyz$ (đơn vị: km), hai trạm phát sóng đặt tại $A(-1;-2;3)$ và $B(1;5;-2)$. Một trạm trung chuyển được đặt tại trung điểm của đoạn thẳng nối $A$ và $B$. Tọa độ của trạm trung chuyển là}
  {$(0;\dfrac{3}{2};\dfrac{1}{2})$.}
  {$(0;\dfrac{7}{2};\dfrac{1}{2})$.}{$(0;\dfrac{1}{2};\dfrac{1}{2})$.}{$(0;-\dfrac{3}{2};\dfrac{1}{2})$.}

\cauTN[id=C2-B8-D11-NB-007,muc=NB,loigiai={Trạm trung chuyển đặt tại trung điểm $I$ của $AB$: $I=(0;-3;\dfrac{3}{2})$.}]{Trong không gian $Oxyz$ (đơn vị: km), hai trạm phát sóng đặt tại $A(-5;-4;-2)$ và $B(5;-2;5)$. Một trạm trung chuyển được đặt tại trung điểm của đoạn thẳng nối $A$ và $B$. Tọa độ của trạm trung chuyển là}
  {$(0;-3;\dfrac{3}{2})$.}
  {$(0;-3;\dfrac{7}{2})$.}{$(\dfrac{3}{2};-3;0)$.}{$(0;3;\dfrac{3}{2})$.}

\cauTN[id=C2-B8-D11-NB-008,muc=NB,loigiai={Trạm trung chuyển đặt tại trung điểm $I$ của $AB$: $I=(-\dfrac{1}{2};-3;-\dfrac{9}{2})$.}]{Trong không gian $Oxyz$ (đơn vị: km), hai trạm phát sóng đặt tại $A(3;-4;-4)$ và $B(-4;-2;-5)$. Một trạm trung chuyển được đặt tại trung điểm của đoạn thẳng nối $A$ và $B$. Tọa độ của trạm trung chuyển là}
  {$(-\dfrac{1}{2};-3;-\dfrac{9}{2})$.}
  {$(-\dfrac{1}{2};3;-\dfrac{9}{2})$.}{$(\dfrac{3}{2};-3;-\dfrac{9}{2})$.}{$(-7;2;-1)$.}

\cauTN[id=C2-B8-D11-NB-009,muc=NB,loigiai={Trạm trung chuyển đặt tại trung điểm $I$ của $AB$: $I=(-\dfrac{1}{2};\dfrac{9}{2};-\dfrac{1}{2})$.}]{Trong không gian $Oxyz$ (đơn vị: km), hai trạm phát sóng đặt tại $A(-3;5;3)$ và $B(2;4;-4)$. Một trạm trung chuyển được đặt tại trung điểm của đoạn thẳng nối $A$ và $B$. Tọa độ của trạm trung chuyển là}
  {$(-\dfrac{1}{2};\dfrac{9}{2};-\dfrac{1}{2})$.}
  {$(-\dfrac{1}{2};-\dfrac{1}{2};\dfrac{9}{2})$.}{$(\dfrac{1}{2};\dfrac{9}{2};-\dfrac{1}{2})$.}{$(-\dfrac{1}{2};\dfrac{11}{2};-\dfrac{1}{2})$.}

\cauTN[id=C2-B8-D11-NB-010,muc=NB,loigiai={Trạm trung chuyển đặt tại trung điểm $I$ của $AB$: $I=(-\dfrac{1}{2};-\dfrac{1}{2};-4)$.}]{Trong không gian $Oxyz$ (đơn vị: km), hai trạm phát sóng đặt tại $A(-1;-5;-3)$ và $B(0;4;-5)$. Một trạm trung chuyển được đặt tại trung điểm của đoạn thẳng nối $A$ và $B$. Tọa độ của trạm trung chuyển là}
  {$(-\dfrac{1}{2};-\dfrac{1}{2};-4)$.}
  {$(-\dfrac{5}{2};-\dfrac{1}{2};-4)$.}{$(-4;-\dfrac{1}{2};-\dfrac{1}{2})$.}{$(-\dfrac{1}{2};-\dfrac{1}{2};4)$.}

\cauTN[id=C2-B8-D11-NB-011,muc=NB,loigiai={Trạm trung chuyển đặt tại trung điểm $I$ của $AB$: $I=(-3;3;2)$.}]{Trong không gian $Oxyz$ (đơn vị: km), hai trạm phát sóng đặt tại $A(-4;1;3)$ và $B(-2;5;1)$. Một trạm trung chuyển được đặt tại trung điểm của đoạn thẳng nối $A$ và $B$. Tọa độ của trạm trung chuyển là}
  {$(-3;3;2)$.}
  {$(-5;3;2)$.}{$(-3;2;3)$.}{$(-3;-3;2)$.}

\cauTN[id=C2-B8-D11-NB-012,muc=NB,loigiai={Trạm trung chuyển đặt tại trung điểm $I$ của $AB$: $I=(4;-\dfrac{1}{2};0)$.}]{Trong không gian $Oxyz$ (đơn vị: km), hai trạm phát sóng đặt tại $A(3;-1;-2)$ và $B(5;0;2)$. Một trạm trung chuyển được đặt tại trung điểm của đoạn thẳng nối $A$ và $B$. Tọa độ của trạm trung chuyển là}
  {$(4;-\dfrac{1}{2};0)$.}
  {$(4;\dfrac{1}{2};0)$.}{$(4;\dfrac{3}{2};0)$.}{$(4;-\dfrac{1}{2};-1)$.}

\cauTN[id=C2-B8-D11-NB-013,muc=NB,loigiai={Khoảng cách giữa hai drone là $AB=\sqrt{(-5)^2+(-5)^2+(-5)^2}=5\sqrt{3}$ (m).}]{Trong không gian $Oxyz$ (đơn vị: m), hai drone cùng lúc ở vị trí $A(2;0;4)$ và $B(-3;-5;-1)$. Khoảng cách giữa hai drone bằng}
  {$5\sqrt{3}$.}
  {$2\sqrt{19}$.}{$\sqrt{74}$.}{$\sqrt{79}$.}

\cauTN[id=C2-B8-D11-NB-014,muc=NB,loigiai={Khoảng cách giữa hai drone là $AB=\sqrt{(1)^2+(3)^2+(-7)^2}=\sqrt{59}$ (m).}]{Trong không gian $Oxyz$ (đơn vị: m), hai drone cùng lúc ở vị trí $A(-1;-1;4)$ và $B(0;2;-3)$. Khoảng cách giữa hai drone bằng}
  {$\sqrt{59}$.}
  {$2\sqrt{15}$.}{$\sqrt{58}$.}{$3\sqrt{7}$.}

\cauTN[id=C2-B8-D11-NB-015,muc=NB,loigiai={Khoảng cách giữa hai drone là $AB=\sqrt{(1)^2+(2)^2+(3)^2}=\sqrt{14}$ (m).}]{Trong không gian $Oxyz$ (đơn vị: m), hai drone cùng lúc ở vị trí $A(-1;0;1)$ và $B(0;2;4)$. Khoảng cách giữa hai drone bằng}
  {$\sqrt{14}$.}
  {$\sqrt{15}$.}{$\sqrt{13}$.}{$3\sqrt{2}$.}

\cauTN[id=C2-B8-D11-NB-016,muc=NB,loigiai={Khoảng cách giữa hai drone là $AB=\sqrt{(-1)^2+(-10)^2+(-7)^2}=5\sqrt{6}$ (m).}]{Trong không gian $Oxyz$ (đơn vị: m), hai drone cùng lúc ở vị trí $A(-1;5;5)$ và $B(-2;-5;-2)$. Khoảng cách giữa hai drone bằng}
  {$5\sqrt{6}$.}
  {$\sqrt{151}$.}{$\sqrt{149}$.}{$\sqrt{154}$.}

\cauTN[id=C2-B8-D11-NB-017,muc=NB,loigiai={Khoảng cách giữa hai drone là $AB=\sqrt{(3)^2+(-2)^2+(8)^2}=\sqrt{77}$ (m).}]{Trong không gian $Oxyz$ (đơn vị: m), hai drone cùng lúc ở vị trí $A(2;2;-5)$ và $B(5;0;3)$. Khoảng cách giữa hai drone bằng}
  {$\sqrt{77}$.}
  {$\sqrt{78}$.}{$2\sqrt{19}$.}{$9$.}

\cauTN[id=C2-B8-D11-NB-018,muc=NB,loigiai={Khoảng cách giữa hai drone là $AB=\sqrt{(1)^2+(-2)^2+(-1)^2}=\sqrt{6}$ (m).}]{Trong không gian $Oxyz$ (đơn vị: m), hai drone cùng lúc ở vị trí $A(0;-2;5)$ và $B(1;-4;4)$. Khoảng cách giữa hai drone bằng}
  {$\sqrt{6}$.}
  {$\sqrt{7}$.}{$\sqrt{5}$.}{$\sqrt{10}$.}

\cauTN[id=C2-B8-D11-NB-019,muc=NB,loigiai={Khoảng cách giữa hai drone là $AB=\sqrt{(-7)^2+(3)^2+(-3)^2}=\sqrt{67}$ (m).}]{Trong không gian $Oxyz$ (đơn vị: m), hai drone cùng lúc ở vị trí $A(2;-4;0)$ và $B(-5;-1;-3)$. Khoảng cách giữa hai drone bằng}
  {$\sqrt{67}$.}
  {$2\sqrt{17}$.}{$\sqrt{66}$.}{$\sqrt{71}$.}

\cauTN[id=C2-B8-D11-NB-020,muc=NB,loigiai={Khoảng cách giữa hai drone là $AB=\sqrt{(4)^2+(-5)^2+(0)^2}=\sqrt{41}$ (m).}]{Trong không gian $Oxyz$ (đơn vị: m), hai drone cùng lúc ở vị trí $A(-3;1;-1)$ và $B(1;-4;-1)$. Khoảng cách giữa hai drone bằng}
  {$\sqrt{41}$.}
  {$\sqrt{42}$.}{$2\sqrt{10}$.}{$3\sqrt{5}$.}

\cauTN[id=C2-B8-D11-NB-021,muc=NB,loigiai={Khoảng cách giữa hai drone là $AB=\sqrt{(2)^2+(-2)^2+(0)^2}=2\sqrt{2}$ (m).}]{Trong không gian $Oxyz$ (đơn vị: m), hai drone cùng lúc ở vị trí $A(-2;-2;2)$ và $B(0;-4;2)$. Khoảng cách giữa hai drone bằng}
  {$2\sqrt{2}$.}
  {$3$.}{$\sqrt{7}$.}{$2\sqrt{3}$.}

\cauTN[id=C2-B8-D11-NB-022,muc=NB,loigiai={Khoảng cách giữa hai drone là $AB=\sqrt{(7)^2+(0)^2+(1)^2}=5\sqrt{2}$ (m).}]{Trong không gian $Oxyz$ (đơn vị: m), hai drone cùng lúc ở vị trí $A(-4;4;-3)$ và $B(3;4;-2)$. Khoảng cách giữa hai drone bằng}
  {$5\sqrt{2}$.}
  {$\sqrt{51}$.}{$7$.}{$3\sqrt{6}$.}

\cauTN[id=C2-B8-D11-NB-023,muc=NB,loigiai={Khoảng cách giữa hai drone là $AB=\sqrt{(-1)^2+(1)^2+(-7)^2}=\sqrt{51}$ (m).}]{Trong không gian $Oxyz$ (đơn vị: m), hai drone cùng lúc ở vị trí $A(2;2;4)$ và $B(1;3;-3)$. Khoảng cách giữa hai drone bằng}
  {$\sqrt{51}$.}
  {$2\sqrt{13}$.}{$5\sqrt{2}$.}{$\sqrt{55}$.}

\cauTN[id=C2-B8-D11-NB-024,muc=NB,loigiai={Khoảng cách giữa hai drone là $AB=\sqrt{(6)^2+(-2)^2+(2)^2}=2\sqrt{11}$ (m).}]{Trong không gian $Oxyz$ (đơn vị: m), hai drone cùng lúc ở vị trí $A(-1;0;-5)$ và $B(5;-2;-3)$. Khoảng cách giữa hai drone bằng}
  {$2\sqrt{11}$.}
  {$3\sqrt{5}$.}{$\sqrt{43}$.}{$4\sqrt{3}$.}

\cauTN[id=C2-B8-D11-NB-025,muc=NB,loigiai={Sau $\dfrac{1}{2}$ giờ, vị trí máy bay là $A+\dfrac{1}{2}\vec{v}=(\dfrac{1}{2};-1;0)$.}]{Trong không gian $Oxyz$ (đơn vị: km), một máy bay lúc xuất phát ở $A(2;-1;1)$ bay thẳng với vận tốc không đổi $\vec{v}=(-3;0;-2)$ (km/h). Sau $\dfrac{1}{2}$ giờ máy bay ở điểm có tọa độ}
  {$(\dfrac{1}{2};-1;0)$.}
  {$(\dfrac{3}{2};-1;0)$.}{$(\dfrac{1}{2};1;0)$.}{$(-\dfrac{1}{2};1;0)$.}

\cauTN[id=C2-B8-D11-NB-026,muc=NB,loigiai={Sau $3$ giờ, vị trí máy bay là $A+3\vec{v}=(-9;7;-10)$.}]{Trong không gian $Oxyz$ (đơn vị: km), một máy bay lúc xuất phát ở $A(-3;4;-1)$ bay thẳng với vận tốc không đổi $\vec{v}=(-2;1;-3)$ (km/h). Sau $3$ giờ máy bay ở điểm có tọa độ}
  {$(-9;7;-10)$.}
  {$(-9;7;-9)$.}{$(-5;5;-4)$.}{$(3;1;8)$.}

\cauTN[id=C2-B8-D11-NB-027,muc=NB,loigiai={Sau $3$ giờ, vị trí máy bay là $A+3\vec{v}=(-13;3;7)$.}]{Trong không gian $Oxyz$ (đơn vị: km), một máy bay lúc xuất phát ở $A(-4;0;-2)$ bay thẳng với vận tốc không đổi $\vec{v}=(-3;1;3)$ (km/h). Sau $3$ giờ máy bay ở điểm có tọa độ}
  {$(-13;3;7)$.}
  {$(-7;1;1)$.}{$(-13;1;7)$.}{$(-13;4;7)$.}

\cauTN[id=C2-B8-D11-NB-028,muc=NB,loigiai={Sau $3$ giờ, vị trí máy bay là $A+3\vec{v}=(-4;-1;-1)$.}]{Trong không gian $Oxyz$ (đơn vị: km), một máy bay lúc xuất phát ở $A(-1;-1;2)$ bay thẳng với vận tốc không đổi $\vec{v}=(-1;0;-1)$ (km/h). Sau $3$ giờ máy bay ở điểm có tọa độ}
  {$(-4;-1;-1)$.}
  {$(-2;-1;-1)$.}{$(-4;0;-1)$.}{$(-4;1;-1)$.}

\cauTN[id=C2-B8-D11-NB-029,muc=NB,loigiai={Sau $\dfrac{1}{2}$ giờ, vị trí máy bay là $A+\dfrac{1}{2}\vec{v}=(-\dfrac{1}{2};-4;\dfrac{1}{2})$.}]{Trong không gian $Oxyz$ (đơn vị: km), một máy bay lúc xuất phát ở $A(-2;-3;1)$ bay thẳng với vận tốc không đổi $\vec{v}=(3;-2;-1)$ (km/h). Sau $\dfrac{1}{2}$ giờ máy bay ở điểm có tọa độ}
  {$(-\dfrac{1}{2};-4;\dfrac{1}{2})$.}
  {$(-\dfrac{1}{2};-2;\dfrac{1}{2})$.}{$(-\dfrac{1}{2};-4;\dfrac{3}{2})$.}{$(-\dfrac{1}{2};-4;-\dfrac{1}{2})$.}

\cauTN[id=C2-B8-D11-NB-030,muc=NB,loigiai={Sau $\dfrac{1}{2}$ giờ, vị trí máy bay là $A+\dfrac{1}{2}\vec{v}=(-\dfrac{1}{2};-\dfrac{5}{2};5)$.}]{Trong không gian $Oxyz$ (đơn vị: km), một máy bay lúc xuất phát ở $A(0;-4;4)$ bay thẳng với vận tốc không đổi $\vec{v}=(-1;3;2)$ (km/h). Sau $\dfrac{1}{2}$ giờ máy bay ở điểm có tọa độ}
  {$(-\dfrac{1}{2};-\dfrac{5}{2};5)$.}
  {$(5;-\dfrac{5}{2};-\dfrac{1}{2})$.}{$(-\dfrac{1}{2};-\dfrac{1}{2};5)$.}{$(-\dfrac{1}{2};-\dfrac{9}{2};5)$.}

\cauTN[id=C2-B8-D11-NB-031,muc=NB,loigiai={Sau $\dfrac{1}{2}$ giờ, vị trí máy bay là $A+\dfrac{1}{2}\vec{v}=(\dfrac{5}{2};2;-\dfrac{7}{2})$.}]{Trong không gian $Oxyz$ (đơn vị: km), một máy bay lúc xuất phát ở $A(3;1;-3)$ bay thẳng với vận tốc không đổi $\vec{v}=(-1;2;-1)$ (km/h). Sau $\dfrac{1}{2}$ giờ máy bay ở điểm có tọa độ}
  {$(\dfrac{5}{2};2;-\dfrac{7}{2})$.}
  {$(\dfrac{7}{2};0;-\dfrac{5}{2})$.}{$(\dfrac{1}{2};2;-\dfrac{7}{2})$.}{$(2;3;-4)$.}

\cauTN[id=C2-B8-D11-NB-032,muc=NB,loigiai={Sau $4$ giờ, vị trí máy bay là $A+4\vec{v}=(-9;-3;9)$.}]{Trong không gian $Oxyz$ (đơn vị: km), một máy bay lúc xuất phát ở $A(-1;-3;1)$ bay thẳng với vận tốc không đổi $\vec{v}=(-2;0;2)$ (km/h). Sau $4$ giờ máy bay ở điểm có tọa độ}
  {$(-9;-3;9)$.}
  {$(-9;-3;10)$.}{$(-9;-3;7)$.}{$(9;-3;9)$.}

\cauTN[id=C2-B8-D11-NB-033,muc=NB,loigiai={Sau $4$ giờ, vị trí máy bay là $A+4\vec{v}=(-11;-12;-8)$.}]{Trong không gian $Oxyz$ (đơn vị: km), một máy bay lúc xuất phát ở $A(1;0;0)$ bay thẳng với vận tốc không đổi $\vec{v}=(-3;-3;-2)$ (km/h). Sau $4$ giờ máy bay ở điểm có tọa độ}
  {$(-11;-12;-8)$.}
  {$(-11;-13;-8)$.}{$(-8;-12;-11)$.}{$(-11;-10;-8)$.}

\cauTN[id=C2-B8-D11-NB-034,muc=NB,loigiai={Sau $\dfrac{1}{2}$ giờ, vị trí máy bay là $A+\dfrac{1}{2}\vec{v}=(-3;\dfrac{5}{2};0)$.}]{Trong không gian $Oxyz$ (đơn vị: km), một máy bay lúc xuất phát ở $A(-3;3;-1)$ bay thẳng với vận tốc không đổi $\vec{v}=(0;-1;2)$ (km/h). Sau $\dfrac{1}{2}$ giờ máy bay ở điểm có tọa độ}
  {$(-3;\dfrac{5}{2};0)$.}
  {$(-3;\dfrac{1}{2};0)$.}{$(3;-\dfrac{5}{2};0)$.}{$(-2;\dfrac{5}{2};0)$.}

\cauTN[id=C2-B8-D11-NB-035,muc=NB,loigiai={Sau $3$ giờ, vị trí máy bay là $A+3\vec{v}=(9;-3;9)$.}]{Trong không gian $Oxyz$ (đơn vị: km), một máy bay lúc xuất phát ở $A(0;0;3)$ bay thẳng với vận tốc không đổi $\vec{v}=(3;-1;2)$ (km/h). Sau $3$ giờ máy bay ở điểm có tọa độ}
  {$(9;-3;9)$.}
  {$(8;-3;9)$.}{$(7;-3;9)$.}{$(9;-3;11)$.}

\cauTN[id=C2-B8-D11-NB-036,muc=NB,loigiai={Sau $2$ giờ, vị trí máy bay là $A+2\vec{v}=(-1;1;4)$.}]{Trong không gian $Oxyz$ (đơn vị: km), một máy bay lúc xuất phát ở $A(-3;-3;2)$ bay thẳng với vận tốc không đổi $\vec{v}=(1;2;1)$ (km/h). Sau $2$ giờ máy bay ở điểm có tọa độ}
  {$(-1;1;4)$.}
  {$(4;1;-1)$.}{$(-2;1;4)$.}{$(0;1;4)$.}

\cauDS[id=C2-B8-D11-NB-037,muc=NB,loigiai={Kiểm tra từng mệnh đề bằng công thức tọa độ. a) Sai. b) Đúng. c) Sai. d) Sai.}]{Xét tính đúng sai của các mệnh đề sau:}{%
  \y{Trong không gian $Oxyz$ (đơn vị: km), một máy bay lúc xuất phát ở $A(3;4;-1)$ bay thẳng với vận tốc không đổi $\vec{v}=(0;3;-1)$ (km/h). Sau $4$ giờ máy bay ở điểm có tọa độ $(3;15;-5)$.}{s}
  \y{Trong không gian $Oxyz$ (đơn vị: km), mặt đất là mặt phẳng $(Oxy)$ và trạm điều khiển đặt tại gốc $O$. Một khinh khí cầu ở vị trí $M(3;-1;7)$. Khoảng cách trên mặt đất từ trạm đến điểm ngay dưới khinh khí cầu bằng $\sqrt{10}$ km.}{d}
  \y{Trong không gian $Oxyz$ (đơn vị: km), hai trạm phát sóng đặt tại $A(2;-4;2)$ và $B(0;-4;-2)$. Trạm trung chuyển đặt tại trung điểm $AB$ có tọa độ $(1;-5;0)$.}{s}
  \y{Trong không gian $Oxyz$ (đơn vị: km), hai trạm phát sóng đặt tại $A(-1;1;1)$ và $B(-4;-3;2)$. Trạm trung chuyển đặt tại trung điểm $AB$ có tọa độ $(-3;-4;1)$.}{s}}

\cauDS[id=C2-B8-D11-NB-038,muc=NB,loigiai={Kiểm tra từng mệnh đề bằng công thức tọa độ. a) Sai. b) Đúng. c) Sai. d) Đúng.}]{Xét tính đúng sai của các mệnh đề sau:}{%
  \y{Trong không gian $Oxyz$ (đơn vị: km), hai trạm phát sóng đặt tại $A(-4;-4;0)$ và $B(1;5;-4)$. Trạm trung chuyển đặt tại trung điểm $AB$ có tọa độ $(-\dfrac{3}{2};\dfrac{3}{2};-2)$.}{s}
  \y{Trong không gian $Oxyz$ (đơn vị: m), hai drone cùng lúc ở vị trí $A(-2;4;1)$ và $B(-4;4;-3)$. Khoảng cách giữa hai drone bằng $2\sqrt{5}$ m.}{d}
  \y{Trong không gian $Oxyz$ (đơn vị: m), hai drone cùng lúc ở vị trí $A(0;-3;-4)$ và $B(5;5;0)$. Khoảng cách giữa hai drone bằng $\sqrt{109}$ m.}{s}
  \y{Trong không gian $Oxyz$ (đơn vị: km), một máy bay lúc xuất phát ở $A(-4;3;0)$ bay thẳng với vận tốc không đổi $\vec{v}=(-3;3;0)$ (km/h). Sau $2$ giờ máy bay ở điểm có tọa độ $(-10;9;0)$.}{d}}

\cauDS[id=C2-B8-D11-NB-039,muc=NB,loigiai={Kiểm tra từng mệnh đề bằng công thức tọa độ. a) Đúng. b) Sai. c) Sai. d) Đúng.}]{Xét tính đúng sai của các mệnh đề sau:}{%
  \y{Trong không gian $Oxyz$ (đơn vị: km), hai trạm phát sóng đặt tại $A(0;0;-1)$ và $B(-4;5;-4)$. Trạm trung chuyển đặt tại trung điểm $AB$ có tọa độ $(-2;\dfrac{5}{2};-\dfrac{5}{2})$.}{d}
  \y{Trong không gian $Oxyz$ (đơn vị: m), hai drone cùng lúc ở vị trí $A(3;-4;1)$ và $B(0;1;-3)$. Khoảng cách giữa hai drone bằng $3\sqrt{6}$ m.}{s}
  \y{Trong không gian $Oxyz$ (đơn vị: km), hai trạm phát sóng đặt tại $A(3;4;3)$ và $B(-4;4;-1)$. Trạm trung chuyển đặt tại trung điểm $AB$ có tọa độ $(-\dfrac{1}{2};3;1)$.}{s}
  \y{Trong không gian $Oxyz$ (đơn vị: km), một máy bay lúc xuất phát ở $A(-3;4;2)$ bay thẳng với vận tốc không đổi $\vec{v}=(0;0;-2)$ (km/h). Sau $3$ giờ máy bay ở điểm có tọa độ $(-3;4;-4)$.}{d}}

\cauDS[id=C2-B8-D11-NB-040,muc=NB,loigiai={Kiểm tra từng mệnh đề bằng công thức tọa độ. a) Sai. b) Đúng. c) Đúng. d) Đúng.}]{Xét tính đúng sai của các mệnh đề sau:}{%
  \y{Trong không gian $Oxyz$ (đơn vị: km), mặt đất là mặt phẳng $(Oxy)$ và trạm điều khiển đặt tại gốc $O$. Một khinh khí cầu ở vị trí $M(-5;-2;6)$. Khoảng cách trên mặt đất từ trạm đến điểm ngay dưới khinh khí cầu bằng $\sqrt{33}$ km.}{s}
  \y{Trong không gian $Oxyz$ (đơn vị: km), hai trạm phát sóng đặt tại $A(0;3;2)$ và $B(-3;-1;-2)$. Trạm trung chuyển đặt tại trung điểm $AB$ có tọa độ $(-\dfrac{3}{2};1;0)$.}{d}
  \y{Trong không gian $Oxyz$ (đơn vị: km), một máy bay lúc xuất phát ở $A(-3;-3;-4)$ bay thẳng với vận tốc không đổi $\vec{v}=(3;-3;-1)$ (km/h). Sau $2$ giờ máy bay ở điểm có tọa độ $(3;-9;-6)$.}{d}
  \y{Trong không gian $Oxyz$ (đơn vị: km), một máy bay lúc xuất phát ở $A(3;-3;2)$ bay thẳng với vận tốc không đổi $\vec{v}=(2;0;-2)$ (km/h). Sau $4$ giờ máy bay ở điểm có tọa độ $(11;-3;-6)$.}{d}}

\cauDS[id=C2-B8-D11-NB-041,muc=NB,loigiai={Kiểm tra từng mệnh đề bằng công thức tọa độ. a) Đúng. b) Đúng. c) Đúng. d) Sai.}]{Xét tính đúng sai của các mệnh đề sau:}{%
  \y{Trong không gian $Oxyz$ (đơn vị: km), hai trạm phát sóng đặt tại $A(3;-1;-2)$ và $B(0;1;-4)$. Trạm trung chuyển đặt tại trung điểm $AB$ có tọa độ $(\dfrac{3}{2};0;-3)$.}{d}
  \y{Trong không gian $Oxyz$ (đơn vị: km), hai trạm phát sóng đặt tại $A(4;2;-5)$ và $B(4;-2;-5)$. Trạm trung chuyển đặt tại trung điểm $AB$ có tọa độ $(4;0;-5)$.}{d}
  \y{Trong không gian $Oxyz$ (đơn vị: m), hai drone cùng lúc ở vị trí $A(-4;3;5)$ và $B(5;2;-3)$. Khoảng cách giữa hai drone bằng $\sqrt{146}$ m.}{d}
  \y{Trong không gian $Oxyz$ (đơn vị: km), mặt đất là mặt phẳng $(Oxy)$ và trạm điều khiển đặt tại gốc $O$. Một khinh khí cầu ở vị trí $M(-3;-6;8)$. Khoảng cách trên mặt đất từ trạm đến điểm ngay dưới khinh khí cầu bằng $7$ km.}{s}}

\cauDS[id=C2-B8-D11-NB-042,muc=NB,loigiai={Kiểm tra từng mệnh đề bằng công thức tọa độ. a) Sai. b) Sai. c) Sai. d) Đúng.}]{Xét tính đúng sai của các mệnh đề sau:}{%
  \y{Trong không gian $Oxyz$ (đơn vị: km), hai trạm phát sóng đặt tại $A(-4;-2;2)$ và $B(-4;-2;-5)$. Trạm trung chuyển đặt tại trung điểm $AB$ có tọa độ $(-4;-2;\dfrac{3}{2})$.}{s}
  \y{Trong không gian $Oxyz$ (đơn vị: km), hai trạm phát sóng đặt tại $A(-1;1;3)$ và $B(2;5;0)$. Trạm trung chuyển đặt tại trung điểm $AB$ có tọa độ $(\dfrac{1}{2};\dfrac{3}{2};3)$.}{s}
  \y{Trong không gian $Oxyz$ (đơn vị: km), mặt đất là mặt phẳng $(Oxy)$ và trạm điều khiển đặt tại gốc $O$. Một khinh khí cầu ở vị trí $M(6;4;4)$. Khoảng cách trên mặt đất từ trạm đến điểm ngay dưới khinh khí cầu bằng $\sqrt{51}$ km.}{s}
  \y{Trong không gian $Oxyz$ (đơn vị: m), hai drone cùng lúc ở vị trí $A(4;0;2)$ và $B(1;1;1)$. Khoảng cách giữa hai drone bằng $\sqrt{11}$ m.}{d}}

\cauDS[id=C2-B8-D11-NB-043,muc=NB,loigiai={Kiểm tra từng mệnh đề bằng công thức tọa độ. a) Đúng. b) Đúng. c) Đúng. d) Sai.}]{Xét tính đúng sai của các mệnh đề sau:}{%
  \y{Trong không gian $Oxyz$ (đơn vị: m), hai drone cùng lúc ở vị trí $A(-2;3;-3)$ và $B(3;2;5)$. Khoảng cách giữa hai drone bằng $3\sqrt{10}$ m.}{d}
  \y{Trong không gian $Oxyz$ (đơn vị: km), hai trạm phát sóng đặt tại $A(5;-1;1)$ và $B(5;-2;-4)$. Trạm trung chuyển đặt tại trung điểm $AB$ có tọa độ $(5;-\dfrac{3}{2};-\dfrac{3}{2})$.}{d}
  \y{Trong không gian $Oxyz$ (đơn vị: km), mặt đất là mặt phẳng $(Oxy)$ và trạm điều khiển đặt tại gốc $O$. Một khinh khí cầu ở vị trí $M(-5;-1;8)$. Khoảng cách trên mặt đất từ trạm đến điểm ngay dưới khinh khí cầu bằng $\sqrt{26}$ km.}{d}
  \y{Trong không gian $Oxyz$ (đơn vị: m), hai drone cùng lúc ở vị trí $A(3;1;5)$ và $B(2;5;-5)$. Khoảng cách giữa hai drone bằng $11$ m.}{s}}

\cauDS[id=C2-B8-D11-NB-044,muc=NB,loigiai={Kiểm tra từng mệnh đề bằng công thức tọa độ. a) Sai. b) Đúng. c) Sai. d) Đúng.}]{Xét tính đúng sai của các mệnh đề sau:}{%
  \y{Trong không gian $Oxyz$ (đơn vị: km), mặt đất là mặt phẳng $(Oxy)$ và trạm điều khiển đặt tại gốc $O$. Một khinh khí cầu ở vị trí $M(3;6;5)$. Khoảng cách trên mặt đất từ trạm đến điểm ngay dưới khinh khí cầu bằng $2\sqrt{11}$ km.}{s}
  \y{Trong không gian $Oxyz$ (đơn vị: m), hai drone cùng lúc ở vị trí $A(-2;2;2)$ và $B(-4;-5;-4)$. Khoảng cách giữa hai drone bằng $\sqrt{89}$ m.}{d}
  \y{Trong không gian $Oxyz$ (đơn vị: km), hai trạm phát sóng đặt tại $A(-2;-3;1)$ và $B(-5;5;-4)$. Trạm trung chuyển đặt tại trung điểm $AB$ có tọa độ $(-\dfrac{7}{2};-1;-\dfrac{3}{2})$.}{s}
  \y{Trong không gian $Oxyz$ (đơn vị: km), một máy bay lúc xuất phát ở $A(-1;-4;-1)$ bay thẳng với vận tốc không đổi $\vec{v}=(3;0;0)$ (km/h). Sau $4$ giờ máy bay ở điểm có tọa độ $(11;-4;-1)$.}{d}}

\cauDS[id=C2-B8-D11-NB-045,muc=NB,loigiai={Kiểm tra từng mệnh đề bằng công thức tọa độ. a) Sai. b) Đúng. c) Sai. d) Đúng.}]{Xét tính đúng sai của các mệnh đề sau:}{%
  \y{Trong không gian $Oxyz$ (đơn vị: m), hai drone cùng lúc ở vị trí $A(-3;-5;0)$ và $B(1;1;-1)$. Khoảng cách giữa hai drone bằng $2\sqrt{13}$ m.}{s}
  \y{Trong không gian $Oxyz$ (đơn vị: km), một máy bay lúc xuất phát ở $A(0;2;-1)$ bay thẳng với vận tốc không đổi $\vec{v}=(3;3;1)$ (km/h). Sau $2$ giờ máy bay ở điểm có tọa độ $(6;8;1)$.}{d}
  \y{Trong không gian $Oxyz$ (đơn vị: km), mặt đất là mặt phẳng $(Oxy)$ và trạm điều khiển đặt tại gốc $O$. Một khinh khí cầu ở vị trí $M(-6;-2;8)$. Khoảng cách trên mặt đất từ trạm đến điểm ngay dưới khinh khí cầu bằng $\sqrt{39}$ km.}{s}
  \y{Trong không gian $Oxyz$ (đơn vị: km), hai trạm phát sóng đặt tại $A(-3;3;4)$ và $B(1;4;-1)$. Trạm trung chuyển đặt tại trung điểm $AB$ có tọa độ $(-1;\dfrac{7}{2};\dfrac{3}{2})$.}{d}}

\cauDS[id=C2-B8-D11-NB-046,muc=NB,loigiai={Kiểm tra từng mệnh đề bằng công thức tọa độ. a) Đúng. b) Sai. c) Sai. d) Sai.}]{Xét tính đúng sai của các mệnh đề sau:}{%
  \y{Trong không gian $Oxyz$ (đơn vị: km), hai trạm phát sóng đặt tại $A(-4;2;4)$ và $B(-5;1;-3)$. Trạm trung chuyển đặt tại trung điểm $AB$ có tọa độ $(-\dfrac{9}{2};\dfrac{3}{2};\dfrac{1}{2})$.}{d}
  \y{Trong không gian $Oxyz$ (đơn vị: km), hai trạm phát sóng đặt tại $A(-1;-2;-5)$ và $B(4;2;-1)$. Trạm trung chuyển đặt tại trung điểm $AB$ có tọa độ $(\dfrac{3}{2};0;-2)$.}{s}
  \y{Trong không gian $Oxyz$ (đơn vị: km), mặt đất là mặt phẳng $(Oxy)$ và trạm điều khiển đặt tại gốc $O$. Một khinh khí cầu ở vị trí $M(6;-2;4)$. Khoảng cách trên mặt đất từ trạm đến điểm ngay dưới khinh khí cầu bằng $2\sqrt{11}$ km.}{s}
  \y{Trong không gian $Oxyz$ (đơn vị: m), hai drone cùng lúc ở vị trí $A(5;4;-1)$ và $B(-3;3;3)$. Khoảng cách giữa hai drone bằng $\sqrt{82}$ m.}{s}}

\cauDS[id=C2-B8-D11-NB-047,muc=NB,loigiai={Kiểm tra từng mệnh đề bằng công thức tọa độ. a) Sai. b) Sai. c) Đúng. d) Đúng.}]{Xét tính đúng sai của các mệnh đề sau:}{%
  \y{Trong không gian $Oxyz$ (đơn vị: km), mặt đất là mặt phẳng $(Oxy)$ và trạm điều khiển đặt tại gốc $O$. Một khinh khí cầu ở vị trí $M(1;-6;3)$. Khoảng cách trên mặt đất từ trạm đến điểm ngay dưới khinh khí cầu bằng $\sqrt{41}$ km.}{s}
  \y{Trong không gian $Oxyz$ (đơn vị: km), mặt đất là mặt phẳng $(Oxy)$ và trạm điều khiển đặt tại gốc $O$. Một khinh khí cầu ở vị trí $M(4;6;3)$. Khoảng cách trên mặt đất từ trạm đến điểm ngay dưới khinh khí cầu bằng $\sqrt{53}$ km.}{s}
  \y{Trong không gian $Oxyz$ (đơn vị: km), hai trạm phát sóng đặt tại $A(1;0;-5)$ và $B(-5;1;2)$. Trạm trung chuyển đặt tại trung điểm $AB$ có tọa độ $(-2;\dfrac{1}{2};-\dfrac{3}{2})$.}{d}
  \y{Trong không gian $Oxyz$ (đơn vị: m), hai drone cùng lúc ở vị trí $A(3;5;-5)$ và $B(-1;-5;-5)$. Khoảng cách giữa hai drone bằng $2\sqrt{29}$ m.}{d}}

\cauDS[id=C2-B8-D11-NB-048,muc=NB,loigiai={Kiểm tra từng mệnh đề bằng công thức tọa độ. a) Đúng. b) Đúng. c) Sai. d) Đúng.}]{Xét tính đúng sai của các mệnh đề sau:}{%
  \y{Trong không gian $Oxyz$ (đơn vị: km), hai trạm phát sóng đặt tại $A(2;3;-2)$ và $B(-2;3;-3)$. Trạm trung chuyển đặt tại trung điểm $AB$ có tọa độ $(0;3;-\dfrac{5}{2})$.}{d}
  \y{Trong không gian $Oxyz$ (đơn vị: km), mặt đất là mặt phẳng $(Oxy)$ và trạm điều khiển đặt tại gốc $O$. Một khinh khí cầu ở vị trí $M(3;6;8)$. Khoảng cách trên mặt đất từ trạm đến điểm ngay dưới khinh khí cầu bằng $3\sqrt{5}$ km.}{d}
  \y{Trong không gian $Oxyz$ (đơn vị: km), một máy bay lúc xuất phát ở $A(0;0;1)$ bay thẳng với vận tốc không đổi $\vec{v}=(1;2;1)$ (km/h). Sau $2$ giờ máy bay ở điểm có tọa độ $(2;5;3)$.}{s}
  \y{Trong không gian $Oxyz$ (đơn vị: m), hai drone cùng lúc ở vị trí $A(5;0;1)$ và $B(4;-4;5)$. Khoảng cách giữa hai drone bằng $\sqrt{33}$ m.}{d}}

\cauTLN[id=C2-B8-D11-NB-049,muc=NB,loigiai={Trạm trung chuyển đặt tại trung điểm $I$ của $AB$: $I=(-5;2;-2)$.}]{Trong không gian $Oxyz$ (đơn vị: km), hai trạm phát sóng đặt tại $A(-5;-1;1)$ và $B(-5;5;-5)$. Một trạm trung chuyển đặt tại trung điểm của $AB$. Tính tổng ba tọa độ của trạm trung chuyển.}{-5}

\cauTLN[id=C2-B8-D11-NB-050,muc=NB,loigiai={Trạm trung chuyển đặt tại trung điểm $I$ của $AB$: $I=(-2;2;1)$.}]{Trong không gian $Oxyz$ (đơn vị: km), hai trạm phát sóng đặt tại $A(-1;3;1)$ và $B(-3;1;1)$. Một trạm trung chuyển đặt tại trung điểm của $AB$. Tính tổng ba tọa độ của trạm trung chuyển.}{1}

\cauTLN[id=C2-B8-D11-NB-051,muc=NB,loigiai={Trạm trung chuyển đặt tại trung điểm $I$ của $AB$: $I=(0;2;2)$.}]{Trong không gian $Oxyz$ (đơn vị: km), hai trạm phát sóng đặt tại $A(-4;-1;0)$ và $B(4;5;4)$. Một trạm trung chuyển đặt tại trung điểm của $AB$. Tính tổng ba tọa độ của trạm trung chuyển.}{4}

\cauTLN[id=C2-B8-D11-NB-052,muc=NB,loigiai={Trạm trung chuyển đặt tại trung điểm $I$ của $AB$: $I=(2;1;3)$.}]{Trong không gian $Oxyz$ (đơn vị: km), hai trạm phát sóng đặt tại $A(1;2;1)$ và $B(3;0;5)$. Một trạm trung chuyển đặt tại trung điểm của $AB$. Tính tổng ba tọa độ của trạm trung chuyển.}{6}

\cauTLN[id=C2-B8-D11-NB-053,muc=NB,loigiai={Khoảng cách giữa hai drone là $AB=\sqrt{(-1)^2+(2)^2+(-3)^2}=\sqrt{14}$ (m).}]{Trong không gian $Oxyz$ (đơn vị: m), hai drone cùng lúc ở vị trí $A(2;-4;-2)$ và $B(1;-2;-5)$. Tính bình phương khoảng cách giữa hai drone (đơn vị: $\text{m}^2$).}{14}

\cauTLN[id=C2-B8-D11-NB-054,muc=NB,loigiai={Khoảng cách giữa hai drone là $AB=\sqrt{(2)^2+(-3)^2+(-3)^2}=\sqrt{22}$ (m).}]{Trong không gian $Oxyz$ (đơn vị: m), hai drone cùng lúc ở vị trí $A(0;0;-2)$ và $B(2;-3;-5)$. Tính bình phương khoảng cách giữa hai drone (đơn vị: $\text{m}^2$).}{22}

\cauTLN[id=C2-B8-D11-NB-055,muc=NB,loigiai={Khoảng cách giữa hai drone là $AB=\sqrt{(-1)^2+(-8)^2+(-4)^2}=9$ (m).}]{Trong không gian $Oxyz$ (đơn vị: m), hai drone cùng lúc ở vị trí $A(-1;4;0)$ và $B(-2;-4;-4)$. Tính bình phương khoảng cách giữa hai drone (đơn vị: $\text{m}^2$).}{81}

\cauTLN[id=C2-B8-D11-NB-056,muc=NB,loigiai={Khoảng cách giữa hai drone là $AB=\sqrt{(-1)^2+(-4)^2+(5)^2}=\sqrt{42}$ (m).}]{Trong không gian $Oxyz$ (đơn vị: m), hai drone cùng lúc ở vị trí $A(0;1;-2)$ và $B(-1;-3;3)$. Tính bình phương khoảng cách giữa hai drone (đơn vị: $\text{m}^2$).}{42}

\cauTLN[id=C2-B8-D11-NB-057,muc=NB,loigiai={Sau $3$ giờ, vị trí máy bay là $A+3\vec{v}=(6;3;-3)$.}]{Trong không gian $Oxyz$ (đơn vị: km), một máy bay lúc xuất phát ở $A(-3;0;0)$ bay thẳng với vận tốc không đổi $\vec{v}=(3;1;-1)$ (km/h). Sau $3$ giờ máy bay ở điểm $M$. Tính $x_M+y_M+z_M$.}{6}

\cauTLN[id=C2-B8-D11-NB-058,muc=NB,loigiai={Sau $4$ giờ, vị trí máy bay là $A+4\vec{v}=(-10;10;-6)$.}]{Trong không gian $Oxyz$ (đơn vị: km), một máy bay lúc xuất phát ở $A(-2;2;-2)$ bay thẳng với vận tốc không đổi $\vec{v}=(-2;2;-1)$ (km/h). Sau $4$ giờ máy bay ở điểm $M$. Tính $x_M+y_M+z_M$.}{-6}

\cauTLN[id=C2-B8-D11-NB-059,muc=NB,loigiai={Sau $4$ giờ, vị trí máy bay là $A+4\vec{v}=(-4;-10;8)$.}]{Trong không gian $Oxyz$ (đơn vị: km), một máy bay lúc xuất phát ở $A(4;-2;-4)$ bay thẳng với vận tốc không đổi $\vec{v}=(-2;-2;3)$ (km/h). Sau $4$ giờ máy bay ở điểm $M$. Tính $x_M+y_M+z_M$.}{-6}

\cauTLN[id=C2-B8-D11-NB-060,muc=NB,loigiai={Sau $3$ giờ, vị trí máy bay là $A+3\vec{v}=(7;7;-10)$.}]{Trong không gian $Oxyz$ (đơn vị: km), một máy bay lúc xuất phát ở $A(-2;4;-1)$ bay thẳng với vận tốc không đổi $\vec{v}=(3;1;-3)$ (km/h). Sau $3$ giờ máy bay ở điểm $M$. Tính $x_M+y_M+z_M$.}{4}

\cauTN[id=C2-B8-D11-TH-001,muc=TH,loigiai={Hình chiếu của $M$ lên mặt đất là $H(6;5;0)$ nên khoảng cách theo mặt đất từ trạm đến điểm ngay dưới khinh khí cầu là $OH=\sqrt{6^2+5^2}=\sqrt{61}$ (km).}]{Trong không gian $Oxyz$ (đơn vị: km), mặt đất là mặt phẳng $(Oxy)$ và trạm điều khiển đặt tại gốc $O$. Một khinh khí cầu ở vị trí $M(6;5;8)$. Khoảng cách trên mặt đất từ trạm điều khiển đến điểm ngay dưới khinh khí cầu bằng}
  {$\sqrt{61}$.}
  {$\sqrt{62}$.}{$2\sqrt{15}$.}{$\sqrt{65}$.}

\cauTN[id=C2-B8-D11-TH-002,muc=TH,loigiai={Hình chiếu của $M$ lên mặt đất là $H(-3;-6;0)$ nên khoảng cách theo mặt đất từ trạm đến điểm ngay dưới khinh khí cầu là $OH=\sqrt{-3^2-6^2}=3\sqrt{5}$ (km).}]{Trong không gian $Oxyz$ (đơn vị: km), mặt đất là mặt phẳng $(Oxy)$ và trạm điều khiển đặt tại gốc $O$. Một khinh khí cầu ở vị trí $M(-3;-6;7)$. Khoảng cách trên mặt đất từ trạm điều khiển đến điểm ngay dưới khinh khí cầu bằng}
  {$3\sqrt{5}$.}
  {$\sqrt{46}$.}{$2\sqrt{11}$.}{$7$.}

\cauTN[id=C2-B8-D11-TH-003,muc=TH,loigiai={Hình chiếu của $M$ lên mặt đất là $H(3;4;0)$ nên khoảng cách theo mặt đất từ trạm đến điểm ngay dưới khinh khí cầu là $OH=\sqrt{3^2+4^2}=5$ (km).}]{Trong không gian $Oxyz$ (đơn vị: km), mặt đất là mặt phẳng $(Oxy)$ và trạm điều khiển đặt tại gốc $O$. Một khinh khí cầu ở vị trí $M(3;4;6)$. Khoảng cách trên mặt đất từ trạm điều khiển đến điểm ngay dưới khinh khí cầu bằng}
  {$5$.}
  {$\sqrt{26}$.}{$2\sqrt{6}$.}{$\sqrt{29}$.}

\cauTN[id=C2-B8-D11-TH-004,muc=TH,loigiai={Hình chiếu của $M$ lên mặt đất là $H(2;-5;0)$ nên khoảng cách theo mặt đất từ trạm đến điểm ngay dưới khinh khí cầu là $OH=\sqrt{2^2-5^2}=\sqrt{29}$ (km).}]{Trong không gian $Oxyz$ (đơn vị: km), mặt đất là mặt phẳng $(Oxy)$ và trạm điều khiển đặt tại gốc $O$. Một khinh khí cầu ở vị trí $M(2;-5;4)$. Khoảng cách trên mặt đất từ trạm điều khiển đến điểm ngay dưới khinh khí cầu bằng}
  {$\sqrt{29}$.}
  {$\sqrt{30}$.}{$2\sqrt{7}$.}{$\sqrt{33}$.}

\cauTN[id=C2-B8-D11-TH-005,muc=TH,loigiai={Hình chiếu của $M$ lên mặt đất là $H(-5;-2;0)$ nên khoảng cách theo mặt đất từ trạm đến điểm ngay dưới khinh khí cầu là $OH=\sqrt{-5^2-2^2}=\sqrt{29}$ (km).}]{Trong không gian $Oxyz$ (đơn vị: km), mặt đất là mặt phẳng $(Oxy)$ và trạm điều khiển đặt tại gốc $O$. Một khinh khí cầu ở vị trí $M(-5;-2;2)$. Khoảng cách trên mặt đất từ trạm điều khiển đến điểm ngay dưới khinh khí cầu bằng}
  {$\sqrt{29}$.}
  {$\sqrt{30}$.}{$2\sqrt{7}$.}{$\sqrt{33}$.}

\cauTN[id=C2-B8-D11-TH-006,muc=TH,loigiai={Hình chiếu của $M$ lên mặt đất là $H(-3;5;0)$ nên khoảng cách theo mặt đất từ trạm đến điểm ngay dưới khinh khí cầu là $OH=\sqrt{-3^2+5^2}=\sqrt{34}$ (km).}]{Trong không gian $Oxyz$ (đơn vị: km), mặt đất là mặt phẳng $(Oxy)$ và trạm điều khiển đặt tại gốc $O$. Một khinh khí cầu ở vị trí $M(-3;5;8)$. Khoảng cách trên mặt đất từ trạm điều khiển đến điểm ngay dưới khinh khí cầu bằng}
  {$\sqrt{34}$.}
  {$\sqrt{35}$.}{$\sqrt{33}$.}{$\sqrt{38}$.}

\cauTN[id=C2-B8-D11-TH-007,muc=TH,loigiai={Hình chiếu của $M$ lên mặt đất là $H(4;-3;0)$ nên khoảng cách theo mặt đất từ trạm đến điểm ngay dưới khinh khí cầu là $OH=\sqrt{4^2-3^2}=5$ (km).}]{Trong không gian $Oxyz$ (đơn vị: km), mặt đất là mặt phẳng $(Oxy)$ và trạm điều khiển đặt tại gốc $O$. Một khinh khí cầu ở vị trí $M(4;-3;1)$. Khoảng cách trên mặt đất từ trạm điều khiển đến điểm ngay dưới khinh khí cầu bằng}
  {$5$.}
  {$\sqrt{26}$.}{$2\sqrt{6}$.}{$\sqrt{29}$.}

\cauTN[id=C2-B8-D11-TH-008,muc=TH,loigiai={Hình chiếu của $M$ lên mặt đất là $H(2;-2;0)$ nên khoảng cách theo mặt đất từ trạm đến điểm ngay dưới khinh khí cầu là $OH=\sqrt{2^2-2^2}=2\sqrt{2}$ (km).}]{Trong không gian $Oxyz$ (đơn vị: km), mặt đất là mặt phẳng $(Oxy)$ và trạm điều khiển đặt tại gốc $O$. Một khinh khí cầu ở vị trí $M(2;-2;1)$. Khoảng cách trên mặt đất từ trạm điều khiển đến điểm ngay dưới khinh khí cầu bằng}
  {$2\sqrt{2}$.}
  {$3$.}{$\sqrt{7}$.}{$2\sqrt{3}$.}

\cauTN[id=C2-B8-D11-TH-009,muc=TH,loigiai={Hình chiếu của $M$ lên mặt đất là $H(1;3;0)$ nên khoảng cách theo mặt đất từ trạm đến điểm ngay dưới khinh khí cầu là $OH=\sqrt{1^2+3^2}=\sqrt{10}$ (km).}]{Trong không gian $Oxyz$ (đơn vị: km), mặt đất là mặt phẳng $(Oxy)$ và trạm điều khiển đặt tại gốc $O$. Một khinh khí cầu ở vị trí $M(1;3;5)$. Khoảng cách trên mặt đất từ trạm điều khiển đến điểm ngay dưới khinh khí cầu bằng}
  {$\sqrt{10}$.}
  {$\sqrt{11}$.}{$3$.}{$\sqrt{14}$.}

\cauTN[id=C2-B8-D11-TH-010,muc=TH,loigiai={Hình chiếu của $M$ lên mặt đất là $H(4;-5;0)$ nên khoảng cách theo mặt đất từ trạm đến điểm ngay dưới khinh khí cầu là $OH=\sqrt{4^2-5^2}=\sqrt{41}$ (km).}]{Trong không gian $Oxyz$ (đơn vị: km), mặt đất là mặt phẳng $(Oxy)$ và trạm điều khiển đặt tại gốc $O$. Một khinh khí cầu ở vị trí $M(4;-5;4)$. Khoảng cách trên mặt đất từ trạm điều khiển đến điểm ngay dưới khinh khí cầu bằng}
  {$\sqrt{41}$.}
  {$\sqrt{42}$.}{$2\sqrt{10}$.}{$3\sqrt{5}$.}

\cauTN[id=C2-B8-D11-TH-011,muc=TH,loigiai={Hình chiếu của $M$ lên mặt đất là $H(0;3;0)$ nên khoảng cách theo mặt đất từ trạm đến điểm ngay dưới khinh khí cầu là $OH=\sqrt{0^2+3^2}=3$ (km).}]{Trong không gian $Oxyz$ (đơn vị: km), mặt đất là mặt phẳng $(Oxy)$ và trạm điều khiển đặt tại gốc $O$. Một khinh khí cầu ở vị trí $M(0;3;2)$. Khoảng cách trên mặt đất từ trạm điều khiển đến điểm ngay dưới khinh khí cầu bằng}
  {$3$.}
  {$\sqrt{10}$.}{$2\sqrt{2}$.}{$\sqrt{13}$.}

\cauTN[id=C2-B8-D11-TH-012,muc=TH,loigiai={Hình chiếu của $M$ lên mặt đất là $H(6;-6;0)$ nên khoảng cách theo mặt đất từ trạm đến điểm ngay dưới khinh khí cầu là $OH=\sqrt{6^2-6^2}=6\sqrt{2}$ (km).}]{Trong không gian $Oxyz$ (đơn vị: km), mặt đất là mặt phẳng $(Oxy)$ và trạm điều khiển đặt tại gốc $O$. Một khinh khí cầu ở vị trí $M(6;-6;7)$. Khoảng cách trên mặt đất từ trạm điều khiển đến điểm ngay dưới khinh khí cầu bằng}
  {$6\sqrt{2}$.}
  {$\sqrt{73}$.}{$\sqrt{71}$.}{$2\sqrt{19}$.}

\cauTN[id=C2-B8-D11-TH-013,muc=TH,loigiai={$AB=\sqrt{(2)^2+(6)^2+(3)^2}=7$ (m).}]{Trong không gian $Oxyz$ (đơn vị: m), một tuyến cáp treo thẳng nối hai trạm tại $A(-3;-1;1)$ và $B(-1;5;4)$. Chiều dài tuyến cáp treo (đơn vị: m) bằng}
  {$7$.}
  {$8$.}{$6$.}{$9$.}

\cauTN[id=C2-B8-D11-TH-014,muc=TH,loigiai={$AB=\sqrt{(3)^2+(0)^2+(4)^2}=5$ (m).}]{Trong không gian $Oxyz$ (đơn vị: m), một tuyến cáp treo thẳng nối hai trạm tại $A(-2;1;-2)$ và $B(1;1;2)$. Chiều dài tuyến cáp treo (đơn vị: m) bằng}
  {$5$.}
  {$6$.}{$4$.}{$7$.}

\cauTN[id=C2-B8-D11-TH-015,muc=TH,loigiai={$AB=\sqrt{(0)^2+(10)^2+(0)^2}=10$ (m).}]{Trong không gian $Oxyz$ (đơn vị: m), một tuyến cáp treo thẳng nối hai trạm tại $A(-2;0;-1)$ và $B(-2;10;-1)$. Chiều dài tuyến cáp treo (đơn vị: m) bằng}
  {$10$.}
  {$11$.}{$9$.}{$12$.}

\cauTN[id=C2-B8-D11-TH-016,muc=TH,loigiai={$AB=\sqrt{(-4)^2+(4)^2+(2)^2}=6$ (m).}]{Trong không gian $Oxyz$ (đơn vị: m), một tuyến cáp treo thẳng nối hai trạm tại $A(-3;1;-3)$ và $B(-7;5;-1)$. Chiều dài tuyến cáp treo (đơn vị: m) bằng}
  {$6$.}
  {$7$.}{$5$.}{$8$.}

\cauTN[id=C2-B8-D11-TH-017,muc=TH,loigiai={$AB=\sqrt{(0)^2+(0)^2+(10)^2}=10$ (m).}]{Trong không gian $Oxyz$ (đơn vị: m), một tuyến cáp treo thẳng nối hai trạm tại $A(-2;2;1)$ và $B(-2;2;11)$. Chiều dài tuyến cáp treo (đơn vị: m) bằng}
  {$10$.}
  {$11$.}{$9$.}{$12$.}

\cauTN[id=C2-B8-D11-TH-018,muc=TH,loigiai={$AB=\sqrt{(4)^2+(6)^2+(12)^2}=14$ (m).}]{Trong không gian $Oxyz$ (đơn vị: m), một tuyến cáp treo thẳng nối hai trạm tại $A(0;1;2)$ và $B(4;7;14)$. Chiều dài tuyến cáp treo (đơn vị: m) bằng}
  {$14$.}
  {$15$.}{$13$.}{$16$.}

\cauTN[id=C2-B8-D11-TH-019,muc=TH,loigiai={$AB=\sqrt{(0)^2+(3)^2+(4)^2}=5$ (m).}]{Trong không gian $Oxyz$ (đơn vị: m), một tuyến cáp treo thẳng nối hai trạm tại $A(1;-2;3)$ và $B(1;1;7)$. Chiều dài tuyến cáp treo (đơn vị: m) bằng}
  {$5$.}
  {$6$.}{$4$.}{$7$.}

\cauTN[id=C2-B8-D11-TH-020,muc=TH,loigiai={$AB=\sqrt{(6)^2+(6)^2+(7)^2}=11$ (m).}]{Trong không gian $Oxyz$ (đơn vị: m), một tuyến cáp treo thẳng nối hai trạm tại $A(1;0;-3)$ và $B(7;6;4)$. Chiều dài tuyến cáp treo (đơn vị: m) bằng}
  {$11$.}
  {$12$.}{$10$.}{$13$.}

\cauTN[id=C2-B8-D11-TH-021,muc=TH,loigiai={$AB=\sqrt{(4)^2+(4)^2+(-2)^2}=6$ (m).}]{Trong không gian $Oxyz$ (đơn vị: m), một tuyến cáp treo thẳng nối hai trạm tại $A(-3;1;1)$ và $B(1;5;-1)$. Chiều dài tuyến cáp treo (đơn vị: m) bằng}
  {$6$.}
  {$7$.}{$5$.}{$8$.}

\cauTN[id=C2-B8-D11-TH-022,muc=TH,loigiai={$AB=\sqrt{(10)^2+(0)^2+(0)^2}=10$ (m).}]{Trong không gian $Oxyz$ (đơn vị: m), một tuyến cáp treo thẳng nối hai trạm tại $A(2;2;-3)$ và $B(12;2;-3)$. Chiều dài tuyến cáp treo (đơn vị: m) bằng}
  {$10$.}
  {$11$.}{$9$.}{$12$.}

\cauTN[id=C2-B8-D11-TH-023,muc=TH,loigiai={$AB=\sqrt{(6)^2+(6)^2+(3)^2}=9$ (m).}]{Trong không gian $Oxyz$ (đơn vị: m), một tuyến cáp treo thẳng nối hai trạm tại $A(3;3;2)$ và $B(9;9;5)$. Chiều dài tuyến cáp treo (đơn vị: m) bằng}
  {$9$.}
  {$10$.}{$8$.}{$11$.}

\cauTN[id=C2-B8-D11-TH-024,muc=TH,loigiai={$AB=\sqrt{(0)^2+(15)^2+(0)^2}=15$ (m).}]{Trong không gian $Oxyz$ (đơn vị: m), một tuyến cáp treo thẳng nối hai trạm tại $A(-1;2;-3)$ và $B(-1;17;-3)$. Chiều dài tuyến cáp treo (đơn vị: m) bằng}
  {$15$.}
  {$16$.}{$14$.}{$17$.}

\cauTN[id=C2-B8-D11-TH-025,muc=TH,loigiai={$AB=12$ km nên thời gian bay là $\dfrac{12}{24}=\dfrac{1}{2}$ giờ.}]{Trong không gian $Oxyz$ (đơn vị: km), một máy bay bay thẳng từ $A(2;1;3)$ đến $B(10;9;-1)$ với tốc độ không đổi $24$ km/h. Thời gian máy bay bay từ $A$ đến $B$ (đơn vị: giờ) là}
  {$\dfrac{1}{2}$.}
  {$2$.}{$\dfrac{7}{2}$.}{$\dfrac{3}{2}$.}

\cauTN[id=C2-B8-D11-TH-026,muc=TH,loigiai={$AB=3$ km nên thời gian bay là $\dfrac{3}{6}=\dfrac{1}{2}$ giờ.}]{Trong không gian $Oxyz$ (đơn vị: km), một máy bay bay thẳng từ $A(0;1;2)$ đến $B(2;3;1)$ với tốc độ không đổi $6$ km/h. Thời gian máy bay bay từ $A$ đến $B$ (đơn vị: giờ) là}
  {$\dfrac{1}{2}$.}
  {$\dfrac{1}{4}$.}{$\dfrac{5}{2}$.}{$-\dfrac{5}{2}$.}

\cauTN[id=C2-B8-D11-TH-027,muc=TH,loigiai={$AB=10$ km nên thời gian bay là $\dfrac{10}{4}=\dfrac{5}{2}$ giờ.}]{Trong không gian $Oxyz$ (đơn vị: km), một máy bay bay thẳng từ $A(2;1;2)$ đến $B(2;9;8)$ với tốc độ không đổi $4$ km/h. Thời gian máy bay bay từ $A$ đến $B$ (đơn vị: giờ) là}
  {$\dfrac{5}{2}$.}
  {$\dfrac{5}{4}$.}{$6$.}{$\dfrac{9}{2}$.}

\cauTN[id=C2-B8-D11-TH-028,muc=TH,loigiai={$AB=5$ km nên thời gian bay là $\dfrac{5}{5}=1$ giờ.}]{Trong không gian $Oxyz$ (đơn vị: km), một máy bay bay thẳng từ $A(0;-1;1)$ đến $B(3;3;1)$ với tốc độ không đổi $5$ km/h. Thời gian máy bay bay từ $A$ đến $B$ (đơn vị: giờ) là}
  {$1$.}
  {$3$.}{$0$.}{$-2$.}

\cauTN[id=C2-B8-D11-TH-029,muc=TH,loigiai={$AB=21$ km nên thời gian bay là $\dfrac{21}{3}=7$ giờ.}]{Trong không gian $Oxyz$ (đơn vị: km), một máy bay bay thẳng từ $A(0;1;2)$ đến $B(6;10;20)$ với tốc độ không đổi $3$ km/h. Thời gian máy bay bay từ $A$ đến $B$ (đơn vị: giờ) là}
  {$7$.}
  {$10$.}{$6$.}{$8$.}

\cauTN[id=C2-B8-D11-TH-030,muc=TH,loigiai={$AB=9$ km nên thời gian bay là $\dfrac{9}{9}=1$ giờ.}]{Trong không gian $Oxyz$ (đơn vị: km), một máy bay bay thẳng từ $A(-3;-1;2)$ đến $B(3;2;8)$ với tốc độ không đổi $9$ km/h. Thời gian máy bay bay từ $A$ đến $B$ (đơn vị: giờ) là}
  {$1$.}
  {$-2$.}{$4$.}{$0$.}

\cauTN[id=C2-B8-D11-TH-031,muc=TH,loigiai={$AB=10$ km nên thời gian bay là $\dfrac{10}{4}=\dfrac{5}{2}$ giờ.}]{Trong không gian $Oxyz$ (đơn vị: km), một máy bay bay thẳng từ $A(-2;-1;-2)$ đến $B(8;-1;-2)$ với tốc độ không đổi $4$ km/h. Thời gian máy bay bay từ $A$ đến $B$ (đơn vị: giờ) là}
  {$\dfrac{5}{2}$.}
  {$\dfrac{3}{2}$.}{$6$.}{$-\dfrac{5}{2}$.}

\cauTN[id=C2-B8-D11-TH-032,muc=TH,loigiai={$AB=12$ km nên thời gian bay là $\dfrac{12}{4}=3$ giờ.}]{Trong không gian $Oxyz$ (đơn vị: km), một máy bay bay thẳng từ $A(-1;3;0)$ đến $B(7;11;-4)$ với tốc độ không đổi $4$ km/h. Thời gian máy bay bay từ $A$ đến $B$ (đơn vị: giờ) là}
  {$3$.}
  {$1$.}{$0$.}{$5$.}

\cauTN[id=C2-B8-D11-TH-033,muc=TH,loigiai={$AB=33$ km nên thời gian bay là $\dfrac{33}{22}=\dfrac{3}{2}$ giờ.}]{Trong không gian $Oxyz$ (đơn vị: km), một máy bay bay thẳng từ $A(-3;-3;2)$ đến $B(15;15;23)$ với tốc độ không đổi $22$ km/h. Thời gian máy bay bay từ $A$ đến $B$ (đơn vị: giờ) là}
  {$\dfrac{3}{2}$.}
  {$\dfrac{1}{2}$.}{$-\dfrac{3}{2}$.}{$4$.}

\cauTN[id=C2-B8-D11-TH-034,muc=TH,loigiai={$AB=5$ km nên thời gian bay là $\dfrac{5}{5}=1$ giờ.}]{Trong không gian $Oxyz$ (đơn vị: km), một máy bay bay thẳng từ $A(-3;1;0)$ đến $B(1;4;0)$ với tốc độ không đổi $5$ km/h. Thời gian máy bay bay từ $A$ đến $B$ (đơn vị: giờ) là}
  {$1$.}
  {$3$.}{$-1$.}{$2$.}

\cauTN[id=C2-B8-D11-TH-035,muc=TH,loigiai={$AB=6$ km nên thời gian bay là $\dfrac{6}{3}=2$ giờ.}]{Trong không gian $Oxyz$ (đơn vị: km), một máy bay bay thẳng từ $A(0;1;3)$ đến $B(4;3;7)$ với tốc độ không đổi $3$ km/h. Thời gian máy bay bay từ $A$ đến $B$ (đơn vị: giờ) là}
  {$2$.}
  {$3$.}{$5$.}{$1$.}

\cauTN[id=C2-B8-D11-TH-036,muc=TH,loigiai={$AB=3$ km nên thời gian bay là $\dfrac{3}{3}=1$ giờ.}]{Trong không gian $Oxyz$ (đơn vị: km), một máy bay bay thẳng từ $A(-2;2;-1)$ đến $B(0;4;0)$ với tốc độ không đổi $3$ km/h. Thời gian máy bay bay từ $A$ đến $B$ (đơn vị: giờ) là}
  {$1$.}
  {$2$.}{$3$.}{$0$.}

\cauTN[id=C2-B8-D11-TH-037,muc=TH,loigiai={$\cos(\vec{a},\vec{b})=\dfrac{0}{\sqrt{5}\sqrt{9}}$ nên góc giữa hai hướng bay bằng $90^\circ$.}]{Tại gốc tọa độ $O$ có một radar. Hai máy bay bay theo hai hướng có vectơ chỉ phương $\vec{a}=(1;0;-2)$ và $\vec{b}=(-2;2;-1)$. Góc giữa hai hướng bay bằng}
  {$90^\circ$.}
  {$60^\circ$.}{$45^\circ$.}{$135^\circ$.}

\cauTN[id=C2-B8-D11-TH-038,muc=TH,loigiai={$\cos(\vec{a},\vec{b})=\dfrac{0}{\sqrt{2}\sqrt{19}}$ nên góc giữa hai hướng bay bằng $90^\circ$.}]{Tại gốc tọa độ $O$ có một radar. Hai máy bay bay theo hai hướng có vectơ chỉ phương $\vec{a}=(-1;1;0)$ và $\vec{b}=(3;3;1)$. Góc giữa hai hướng bay bằng}
  {$90^\circ$.}
  {$60^\circ$.}{$120^\circ$.}{$30^\circ$.}

\cauTN[id=C2-B8-D11-TH-039,muc=TH,loigiai={$\cos(\vec{a},\vec{b})=\dfrac{0}{\sqrt{27}\sqrt{14}}$ nên góc giữa hai hướng bay bằng $90^\circ$.}]{Tại gốc tọa độ $O$ có một radar. Hai máy bay bay theo hai hướng có vectơ chỉ phương $\vec{a}=(3;3;3)$ và $\vec{b}=(3;-1;-2)$. Góc giữa hai hướng bay bằng}
  {$90^\circ$.}
  {$45^\circ$.}{$30^\circ$.}{$60^\circ$.}

\cauTN[id=C2-B8-D11-TH-040,muc=TH,loigiai={$\cos(\vec{a},\vec{b})=\dfrac{0}{\sqrt{9}\sqrt{18}}$ nên góc giữa hai hướng bay bằng $90^\circ$.}]{Tại gốc tọa độ $O$ có một radar. Hai máy bay bay theo hai hướng có vectơ chỉ phương $\vec{a}=(0;0;-3)$ và $\vec{b}=(3;3;0)$. Góc giữa hai hướng bay bằng}
  {$90^\circ$.}
  {$45^\circ$.}{$60^\circ$.}{$30^\circ$.}

\cauTN[id=C2-B8-D11-TH-041,muc=TH,loigiai={$\cos(\vec{a},\vec{b})=\dfrac{0}{\sqrt{22}\sqrt{18}}$ nên góc giữa hai hướng bay bằng $90^\circ$.}]{Tại gốc tọa độ $O$ có một radar. Hai máy bay bay theo hai hướng có vectơ chỉ phương $\vec{a}=(-2;3;3)$ và $\vec{b}=(0;3;-3)$. Góc giữa hai hướng bay bằng}
  {$90^\circ$.}
  {$150^\circ$.}{$30^\circ$.}{$135^\circ$.}

\cauTN[id=C2-B8-D11-TH-042,muc=TH,loigiai={$\cos(\vec{a},\vec{b})=\dfrac{0}{\sqrt{19}\sqrt{18}}$ nên góc giữa hai hướng bay bằng $90^\circ$.}]{Tại gốc tọa độ $O$ có một radar. Hai máy bay bay theo hai hướng có vectơ chỉ phương $\vec{a}=(-3;-3;1)$ và $\vec{b}=(3;-3;0)$. Góc giữa hai hướng bay bằng}
  {$90^\circ$.}
  {$30^\circ$.}{$60^\circ$.}{$135^\circ$.}

\cauTN[id=C2-B8-D11-TH-043,muc=TH,loigiai={$\cos(\vec{a},\vec{b})=\dfrac{0}{\sqrt{9}\sqrt{9}}$ nên góc giữa hai hướng bay bằng $90^\circ$.}]{Tại gốc tọa độ $O$ có một radar. Hai máy bay bay theo hai hướng có vectơ chỉ phương $\vec{a}=(1;-2;-2)$ và $\vec{b}=(2;-1;2)$. Góc giữa hai hướng bay bằng}
  {$90^\circ$.}
  {$150^\circ$.}{$45^\circ$.}{$120^\circ$.}

\cauTN[id=C2-B8-D11-TH-044,muc=TH,loigiai={$\cos(\vec{a},\vec{b})=\dfrac{0}{\sqrt{5}\sqrt{4}}$ nên góc giữa hai hướng bay bằng $90^\circ$.}]{Tại gốc tọa độ $O$ có một radar. Hai máy bay bay theo hai hướng có vectơ chỉ phương $\vec{a}=(-2;-1;0)$ và $\vec{b}=(0;0;2)$. Góc giữa hai hướng bay bằng}
  {$90^\circ$.}
  {$60^\circ$.}{$135^\circ$.}{$120^\circ$.}

\cauTN[id=C2-B8-D11-TH-045,muc=TH,loigiai={$\cos(\vec{a},\vec{b})=\dfrac{6}{\sqrt{8}\sqrt{6}}$ nên góc giữa hai hướng bay bằng $30^\circ$.}]{Tại gốc tọa độ $O$ có một radar. Hai máy bay bay theo hai hướng có vectơ chỉ phương $\vec{a}=(-2;0;-2)$ và $\vec{b}=(-1;1;-2)$. Góc giữa hai hướng bay bằng}
  {$30^\circ$.}
  {$90^\circ$.}{$150^\circ$.}{$60^\circ$.}

\cauTN[id=C2-B8-D11-TH-046,muc=TH,loigiai={$\cos(\vec{a},\vec{b})=\dfrac{0}{\sqrt{13}\sqrt{17}}$ nên góc giữa hai hướng bay bằng $90^\circ$.}]{Tại gốc tọa độ $O$ có một radar. Hai máy bay bay theo hai hướng có vectơ chỉ phương $\vec{a}=(-2;0;3)$ và $\vec{b}=(3;2;2)$. Góc giữa hai hướng bay bằng}
  {$90^\circ$.}
  {$45^\circ$.}{$30^\circ$.}{$135^\circ$.}

\cauTN[id=C2-B8-D11-TH-047,muc=TH,loigiai={$\cos(\vec{a},\vec{b})=\dfrac{0}{\sqrt{6}\sqrt{14}}$ nên góc giữa hai hướng bay bằng $90^\circ$.}]{Tại gốc tọa độ $O$ có một radar. Hai máy bay bay theo hai hướng có vectơ chỉ phương $\vec{a}=(-2;1;1)$ và $\vec{b}=(2;1;3)$. Góc giữa hai hướng bay bằng}
  {$90^\circ$.}
  {$60^\circ$.}{$120^\circ$.}{$30^\circ$.}

\cauTN[id=C2-B8-D11-TH-048,muc=TH,loigiai={$\cos(\vec{a},\vec{b})=\dfrac{0}{\sqrt{19}\sqrt{22}}$ nên góc giữa hai hướng bay bằng $90^\circ$.}]{Tại gốc tọa độ $O$ có một radar. Hai máy bay bay theo hai hướng có vectơ chỉ phương $\vec{a}=(3;1;3)$ và $\vec{b}=(3;-3;-2)$. Góc giữa hai hướng bay bằng}
  {$90^\circ$.}
  {$60^\circ$.}{$150^\circ$.}{$120^\circ$.}

\cauDS[id=C2-B8-D11-TH-049,muc=TH,loigiai={Kiểm tra từng mệnh đề bằng công thức tọa độ. a) Đúng. b) Đúng. c) Đúng. d) Sai.}]{Xét tính đúng sai của các mệnh đề sau:}{%
  \y{Trong không gian $Oxyz$ (đơn vị: m), một tuyến cáp treo thẳng nối hai trạm tại $A(0;-3;1)$ và $B(1;-5;3)$. Chiều dài tuyến cáp treo bằng $3$ m.}{d}
  \y{Trong không gian $Oxyz$ (đơn vị: km), một máy bay bay thẳng từ $A(1;3;-2)$ đến $B(5;1;2)$ với tốc độ không đổi $6$ km/h. Thời gian bay từ $A$ đến $B$ là $1$ giờ.}{d}
  \y{Trong không gian $Oxyz$ (đơn vị: m), một tuyến cáp treo thẳng nối hai trạm tại $A(-3;3;-2)$ và $B(9;15;12)$. Chiều dài tuyến cáp treo bằng $22$ m.}{d}
  \y{Trong không gian $Oxyz$ (đơn vị: km), mặt đất là mặt phẳng $(Oxy)$ và trạm điều khiển đặt tại gốc $O$. Một khinh khí cầu ở vị trí $M(-3;2;5)$. Khoảng cách trên mặt đất từ trạm đến điểm ngay dưới khinh khí cầu bằng $2\sqrt{3}$ km.}{s}}

\cauDS[id=C2-B8-D11-TH-050,muc=TH,loigiai={Kiểm tra từng mệnh đề bằng công thức tọa độ. a) Sai. b) Đúng. c) Đúng. d) Sai.}]{Xét tính đúng sai của các mệnh đề sau:}{%
  \y{Trong không gian $Oxyz$ (đơn vị: km), mặt đất là mặt phẳng $(Oxy)$ và trạm điều khiển đặt tại gốc $O$. Một khinh khí cầu ở vị trí $M(5;6;6)$. Khoảng cách trên mặt đất từ trạm đến điểm ngay dưới khinh khí cầu bằng $2\sqrt{15}$ km.}{s}
  \y{Trong không gian $Oxyz$ (đơn vị: km), mặt đất là mặt phẳng $(Oxy)$ và trạm điều khiển đặt tại gốc $O$. Một khinh khí cầu ở vị trí $M(-5;3;1)$. Khoảng cách trên mặt đất từ trạm đến điểm ngay dưới khinh khí cầu bằng $\sqrt{34}$ km.}{d}
  \y{Trong không gian $Oxyz$ (đơn vị: km), một máy bay bay thẳng từ $A(2;1;1)$ đến $B(11;7;19)$ với tốc độ không đổi $21$ km/h. Thời gian bay từ $A$ đến $B$ là $1$ giờ.}{d}
  \y{Trong không gian $Oxyz$ (đơn vị: m), một tuyến cáp treo thẳng nối hai trạm tại $A(0;3;-2)$ và $B(2;2;0)$. Chiều dài tuyến cáp treo bằng $2$ m.}{s}}

\cauDS[id=C2-B8-D11-TH-051,muc=TH,loigiai={Kiểm tra từng mệnh đề bằng công thức tọa độ. a) Sai. b) Sai. c) Sai. d) Đúng.}]{Xét tính đúng sai của các mệnh đề sau:}{%
  \y{Tại gốc tọa độ $O$ có một radar. Hai máy bay bay theo hai hướng có vectơ chỉ phương $\vec{a}=(2;-1;2)$ và $\vec{b}=(3;2;-2)$. Góc giữa hai hướng bay bằng $30^\circ$.}{s}
  \y{Trong không gian $Oxyz$ (đơn vị: km), mặt đất là mặt phẳng $(Oxy)$ và trạm điều khiển đặt tại gốc $O$. Một khinh khí cầu ở vị trí $M(-6;-3;4)$. Khoảng cách trên mặt đất từ trạm đến điểm ngay dưới khinh khí cầu bằng $\sqrt{46}$ km.}{s}
  \y{Trong không gian $Oxyz$ (đơn vị: km), một máy bay bay thẳng từ $A(1;-2;1)$ đến $B(10;10;1)$ với tốc độ không đổi $5$ km/h. Thời gian bay từ $A$ đến $B$ là $5$ giờ.}{s}
  \y{Tại gốc tọa độ $O$ có một radar. Hai máy bay bay theo hai hướng có vectơ chỉ phương $\vec{a}=(-1;-1;2)$ và $\vec{b}=(0;2;-2)$. Góc giữa hai hướng bay bằng $150^\circ$.}{d}}

\cauDS[id=C2-B8-D11-TH-052,muc=TH,loigiai={Kiểm tra từng mệnh đề bằng công thức tọa độ. a) Sai. b) Đúng. c) Sai. d) Đúng.}]{Xét tính đúng sai của các mệnh đề sau:}{%
  \y{Tại gốc tọa độ $O$ có một radar. Hai máy bay bay theo hai hướng có vectơ chỉ phương $\vec{a}=(3;-1;2)$ và $\vec{b}=(1;3;0)$. Góc giữa hai hướng bay bằng $150^\circ$.}{s}
  \y{Trong không gian $Oxyz$ (đơn vị: km), một máy bay bay thẳng từ $A(3;2;2)$ đến $B(21;20;23)$ với tốc độ không đổi $6$ km/h. Thời gian bay từ $A$ đến $B$ là $\dfrac{11}{2}$ giờ.}{d}
  \y{Trong không gian $Oxyz$ (đơn vị: m), một tuyến cáp treo thẳng nối hai trạm tại $A(-3;-3;-1)$ và $B(15;15;20)$. Chiều dài tuyến cáp treo bằng $35$ m.}{s}
  \y{Trong không gian $Oxyz$ (đơn vị: km), mặt đất là mặt phẳng $(Oxy)$ và trạm điều khiển đặt tại gốc $O$. Một khinh khí cầu ở vị trí $M(5;-1;3)$. Khoảng cách trên mặt đất từ trạm đến điểm ngay dưới khinh khí cầu bằng $\sqrt{26}$ km.}{d}}

\cauDS[id=C2-B8-D11-TH-053,muc=TH,loigiai={Kiểm tra từng mệnh đề bằng công thức tọa độ. a) Sai. b) Sai. c) Đúng. d) Đúng.}]{Xét tính đúng sai của các mệnh đề sau:}{%
  \y{Trong không gian $Oxyz$ (đơn vị: m), một tuyến cáp treo thẳng nối hai trạm tại $A(0;1;-3)$ và $B(5;1;-3)$. Chiều dài tuyến cáp treo bằng $7$ m.}{s}
  \y{Tại gốc tọa độ $O$ có một radar. Hai máy bay bay theo hai hướng có vectơ chỉ phương $\vec{a}=(-1;0;-1)$ và $\vec{b}=(3;-1;-3)$. Góc giữa hai hướng bay bằng $120^\circ$.}{s}
  \y{Trong không gian $Oxyz$ (đơn vị: km), một máy bay bay thẳng từ $A(2;3;2)$ đến $B(4;7;6)$ với tốc độ không đổi $6$ km/h. Thời gian bay từ $A$ đến $B$ là $1$ giờ.}{d}
  \y{Tại gốc tọa độ $O$ có một radar. Hai máy bay bay theo hai hướng có vectơ chỉ phương $\vec{a}=(0;0;3)$ và $\vec{b}=(0;-3;3)$. Góc giữa hai hướng bay bằng $45^\circ$.}{d}}

\cauDS[id=C2-B8-D11-TH-054,muc=TH,loigiai={Kiểm tra từng mệnh đề bằng công thức tọa độ. a) Sai. b) Sai. c) Đúng. d) Sai.}]{Xét tính đúng sai của các mệnh đề sau:}{%
  \y{Tại gốc tọa độ $O$ có một radar. Hai máy bay bay theo hai hướng có vectơ chỉ phương $\vec{a}=(-1;-3;0)$ và $\vec{b}=(3;-1;1)$. Góc giữa hai hướng bay bằng $135^\circ$.}{s}
  \y{Trong không gian $Oxyz$ (đơn vị: km), một máy bay bay thẳng từ $A(-1;-2;0)$ đến $B(-1;18;0)$ với tốc độ không đổi $4$ km/h. Thời gian bay từ $A$ đến $B$ là $4$ giờ.}{s}
  \y{Trong không gian $Oxyz$ (đơn vị: m), một tuyến cáp treo thẳng nối hai trạm tại $A(-1;1;-2)$ và $B(3;1;1)$. Chiều dài tuyến cáp treo bằng $5$ m.}{d}
  \y{Trong không gian $Oxyz$ (đơn vị: km), mặt đất là mặt phẳng $(Oxy)$ và trạm điều khiển đặt tại gốc $O$. Một khinh khí cầu ở vị trí $M(-2;-2;4)$. Khoảng cách trên mặt đất từ trạm đến điểm ngay dưới khinh khí cầu bằng $\sqrt{7}$ km.}{s}}

\cauDS[id=C2-B8-D11-TH-055,muc=TH,loigiai={Kiểm tra từng mệnh đề bằng công thức tọa độ. a) Đúng. b) Đúng. c) Sai. d) Sai.}]{Xét tính đúng sai của các mệnh đề sau:}{%
  \y{Tại gốc tọa độ $O$ có một radar. Hai máy bay bay theo hai hướng có vectơ chỉ phương $\vec{a}=(2;2;3)$ và $\vec{b}=(3;-3;0)$. Góc giữa hai hướng bay bằng $90^\circ$.}{d}
  \y{Trong không gian $Oxyz$ (đơn vị: km), mặt đất là mặt phẳng $(Oxy)$ và trạm điều khiển đặt tại gốc $O$. Một khinh khí cầu ở vị trí $M(6;1;2)$. Khoảng cách trên mặt đất từ trạm đến điểm ngay dưới khinh khí cầu bằng $\sqrt{37}$ km.}{d}
  \y{Trong không gian $Oxyz$ (đơn vị: m), một tuyến cáp treo thẳng nối hai trạm tại $A(-2;-2;-2)$ và $B(0;2;2)$. Chiều dài tuyến cáp treo bằng $8$ m.}{s}
  \y{Trong không gian $Oxyz$ (đơn vị: km), một máy bay bay thẳng từ $A(-2;0;2)$ đến $B(-2;9;14)$ với tốc độ không đổi $5$ km/h. Thời gian bay từ $A$ đến $B$ là $6$ giờ.}{s}}

\cauDS[id=C2-B8-D11-TH-056,muc=TH,loigiai={Kiểm tra từng mệnh đề bằng công thức tọa độ. a) Sai. b) Sai. c) Đúng. d) Sai.}]{Xét tính đúng sai của các mệnh đề sau:}{%
  \y{Tại gốc tọa độ $O$ có một radar. Hai máy bay bay theo hai hướng có vectơ chỉ phương $\vec{a}=(1;0;-1)$ và $\vec{b}=(0;-2;2)$. Góc giữa hai hướng bay bằng $30^\circ$.}{s}
  \y{Trong không gian $Oxyz$ (đơn vị: km), mặt đất là mặt phẳng $(Oxy)$ và trạm điều khiển đặt tại gốc $O$. Một khinh khí cầu ở vị trí $M(2;0;7)$. Khoảng cách trên mặt đất từ trạm đến điểm ngay dưới khinh khí cầu bằng $\sqrt{5}$ km.}{s}
  \y{Trong không gian $Oxyz$ (đơn vị: km), một máy bay bay thẳng từ $A(-2;0;-2)$ đến $B(1;0;2)$ với tốc độ không đổi $10$ km/h. Thời gian bay từ $A$ đến $B$ là $\dfrac{1}{2}$ giờ.}{d}
  \y{Tại gốc tọa độ $O$ có một radar. Hai máy bay bay theo hai hướng có vectơ chỉ phương $\vec{a}=(-1;-1;2)$ và $\vec{b}=(-3;1;-1)$. Góc giữa hai hướng bay bằng $45^\circ$.}{s}}

\cauDS[id=C2-B8-D11-TH-057,muc=TH,loigiai={Kiểm tra từng mệnh đề bằng công thức tọa độ. a) Đúng. b) Sai. c) Đúng. d) Đúng.}]{Xét tính đúng sai của các mệnh đề sau:}{%
  \y{Tại gốc tọa độ $O$ có một radar. Hai máy bay bay theo hai hướng có vectơ chỉ phương $\vec{a}=(0;2;2)$ và $\vec{b}=(0;-3;0)$. Góc giữa hai hướng bay bằng $135^\circ$.}{d}
  \y{Trong không gian $Oxyz$ (đơn vị: m), một tuyến cáp treo thẳng nối hai trạm tại $A(3;-1;0)$ và $B(9;2;6)$. Chiều dài tuyến cáp treo bằng $11$ m.}{s}
  \y{Trong không gian $Oxyz$ (đơn vị: km), mặt đất là mặt phẳng $(Oxy)$ và trạm điều khiển đặt tại gốc $O$. Một khinh khí cầu ở vị trí $M(5;-5;1)$. Khoảng cách trên mặt đất từ trạm đến điểm ngay dưới khinh khí cầu bằng $5\sqrt{2}$ km.}{d}
  \y{Tại gốc tọa độ $O$ có một radar. Hai máy bay bay theo hai hướng có vectơ chỉ phương $\vec{a}=(-1;-3;3)$ và $\vec{b}=(-3;2;1)$. Góc giữa hai hướng bay bằng $90^\circ$.}{d}}

\cauDS[id=C2-B8-D11-TH-058,muc=TH,loigiai={Kiểm tra từng mệnh đề bằng công thức tọa độ. a) Sai. b) Đúng. c) Đúng. d) Đúng.}]{Xét tính đúng sai của các mệnh đề sau:}{%
  \y{Trong không gian $Oxyz$ (đơn vị: km), một máy bay bay thẳng từ $A(0;0;-3)$ đến $B(6;6;-6)$ với tốc độ không đổi $18$ km/h. Thời gian bay từ $A$ đến $B$ là $\dfrac{3}{2}$ giờ.}{s}
  \y{Trong không gian $Oxyz$ (đơn vị: km), mặt đất là mặt phẳng $(Oxy)$ và trạm điều khiển đặt tại gốc $O$. Một khinh khí cầu ở vị trí $M(-5;-5;7)$. Khoảng cách trên mặt đất từ trạm đến điểm ngay dưới khinh khí cầu bằng $5\sqrt{2}$ km.}{d}
  \y{Trong không gian $Oxyz$ (đơn vị: m), một tuyến cáp treo thẳng nối hai trạm tại $A(-3;2;-2)$ và $B(-7;6;0)$. Chiều dài tuyến cáp treo bằng $6$ m.}{d}
  \y{Trong không gian $Oxyz$ (đơn vị: km), một máy bay bay thẳng từ $A(1;-2;1)$ đến $B(1;7;13)$ với tốc độ không đổi $6$ km/h. Thời gian bay từ $A$ đến $B$ là $\dfrac{5}{2}$ giờ.}{d}}

\cauDS[id=C2-B8-D11-TH-059,muc=TH,loigiai={Kiểm tra từng mệnh đề bằng công thức tọa độ. a) Sai. b) Sai. c) Đúng. d) Sai.}]{Xét tính đúng sai của các mệnh đề sau:}{%
  \y{Trong không gian $Oxyz$ (đơn vị: km), một máy bay bay thẳng từ $A(-2;-2;-1)$ đến $B(0;-3;1)$ với tốc độ không đổi $6$ km/h. Thời gian bay từ $A$ đến $B$ là $-\dfrac{3}{2}$ giờ.}{s}
  \y{Trong không gian $Oxyz$ (đơn vị: m), một tuyến cáp treo thẳng nối hai trạm tại $A(-2;1;-1)$ và $B(10;7;3)$. Chiều dài tuyến cáp treo bằng $13$ m.}{s}
  \y{Trong không gian $Oxyz$ (đơn vị: km), mặt đất là mặt phẳng $(Oxy)$ và trạm điều khiển đặt tại gốc $O$. Một khinh khí cầu ở vị trí $M(-5;3;1)$. Khoảng cách trên mặt đất từ trạm đến điểm ngay dưới khinh khí cầu bằng $\sqrt{34}$ km.}{d}
  \y{Trong không gian $Oxyz$ (đơn vị: km), một máy bay bay thẳng từ $A(-3;1;2)$ đến $B(-3;5;5)$ với tốc độ không đổi $10$ km/h. Thời gian bay từ $A$ đến $B$ là $\dfrac{5}{2}$ giờ.}{s}}

\cauDS[id=C2-B8-D11-TH-060,muc=TH,loigiai={Kiểm tra từng mệnh đề bằng công thức tọa độ. a) Sai. b) Đúng. c) Sai. d) Đúng.}]{Xét tính đúng sai của các mệnh đề sau:}{%
  \y{Tại gốc tọa độ $O$ có một radar. Hai máy bay bay theo hai hướng có vectơ chỉ phương $\vec{a}=(1;-3;-2)$ và $\vec{b}=(2;0;1)$. Góc giữa hai hướng bay bằng $45^\circ$.}{s}
  \y{Trong không gian $Oxyz$ (đơn vị: km), một máy bay bay thẳng từ $A(3;-3;1)$ đến $B(9;3;8)$ với tốc độ không đổi $11$ km/h. Thời gian bay từ $A$ đến $B$ là $1$ giờ.}{d}
  \y{Trong không gian $Oxyz$ (đơn vị: km), mặt đất là mặt phẳng $(Oxy)$ và trạm điều khiển đặt tại gốc $O$. Một khinh khí cầu ở vị trí $M(-4;-3;8)$. Khoảng cách trên mặt đất từ trạm đến điểm ngay dưới khinh khí cầu bằng $2\sqrt{6}$ km.}{s}
  \y{Trong không gian $Oxyz$ (đơn vị: km), mặt đất là mặt phẳng $(Oxy)$ và trạm điều khiển đặt tại gốc $O$. Một khinh khí cầu ở vị trí $M(-1;2;7)$. Khoảng cách trên mặt đất từ trạm đến điểm ngay dưới khinh khí cầu bằng $\sqrt{5}$ km.}{d}}

\cauDS[id=C2-B8-D11-TH-061,muc=TH,loigiai={Kiểm tra từng mệnh đề bằng công thức tọa độ. a) Sai. b) Đúng. c) Sai. d) Đúng.}]{Xét tính đúng sai của các mệnh đề sau:}{%
  \y{Trong không gian $Oxyz$ (đơn vị: km), mặt đất là mặt phẳng $(Oxy)$ và trạm điều khiển đặt tại gốc $O$. Một khinh khí cầu ở vị trí $M(5;1;5)$. Khoảng cách trên mặt đất từ trạm đến điểm ngay dưới khinh khí cầu bằng $3\sqrt{3}$ km.}{s}
  \y{Trong không gian $Oxyz$ (đơn vị: m), một tuyến cáp treo thẳng nối hai trạm tại $A(1;2;1)$ và $B(6;2;1)$. Chiều dài tuyến cáp treo bằng $5$ m.}{d}
  \y{Tại gốc tọa độ $O$ có một radar. Hai máy bay bay theo hai hướng có vectơ chỉ phương $\vec{a}=(-3;0;-2)$ và $\vec{b}=(0;3;0)$. Góc giữa hai hướng bay bằng $45^\circ$.}{s}
  \y{Tại gốc tọa độ $O$ có một radar. Hai máy bay bay theo hai hướng có vectơ chỉ phương $\vec{a}=(-2;0;0)$ và $\vec{b}=(0;3;1)$. Góc giữa hai hướng bay bằng $90^\circ$.}{d}}

\cauDS[id=C2-B8-D11-TH-062,muc=TH,loigiai={Kiểm tra từng mệnh đề bằng công thức tọa độ. a) Đúng. b) Đúng. c) Sai. d) Sai.}]{Xét tính đúng sai của các mệnh đề sau:}{%
  \y{Trong không gian $Oxyz$ (đơn vị: km), một máy bay bay thẳng từ $A(-1;-1;-1)$ đến $B(-5;3;1)$ với tốc độ không đổi $6$ km/h. Thời gian bay từ $A$ đến $B$ là $1$ giờ.}{d}
  \y{Tại gốc tọa độ $O$ có một radar. Hai máy bay bay theo hai hướng có vectơ chỉ phương $\vec{a}=(2;-2;-3)$ và $\vec{b}=(-3;-3;0)$. Góc giữa hai hướng bay bằng $90^\circ$.}{d}
  \y{Trong không gian $Oxyz$ (đơn vị: m), một tuyến cáp treo thẳng nối hai trạm tại $A(1;1;-3)$ và $B(7;9;-3)$. Chiều dài tuyến cáp treo bằng $9$ m.}{s}
  \y{Trong không gian $Oxyz$ (đơn vị: km), một máy bay bay thẳng từ $A(0;-2;-3)$ đến $B(4;6;5)$ với tốc độ không đổi $6$ km/h. Thời gian bay từ $A$ đến $B$ là $-1$ giờ.}{s}}

\cauDS[id=C2-B8-D11-TH-063,muc=TH,loigiai={Kiểm tra từng mệnh đề bằng công thức tọa độ. a) Sai. b) Đúng. c) Đúng. d) Đúng.}]{Xét tính đúng sai của các mệnh đề sau:}{%
  \y{Tại gốc tọa độ $O$ có một radar. Hai máy bay bay theo hai hướng có vectơ chỉ phương $\vec{a}=(3;-1;0)$ và $\vec{b}=(1;3;-3)$. Góc giữa hai hướng bay bằng $60^\circ$.}{s}
  \y{Trong không gian $Oxyz$ (đơn vị: m), một tuyến cáp treo thẳng nối hai trạm tại $A(-1;-1;-2)$ và $B(17;17;19)$. Chiều dài tuyến cáp treo bằng $33$ m.}{d}
  \y{Trong không gian $Oxyz$ (đơn vị: km), mặt đất là mặt phẳng $(Oxy)$ và trạm điều khiển đặt tại gốc $O$. Một khinh khí cầu ở vị trí $M(-4;6;3)$. Khoảng cách trên mặt đất từ trạm đến điểm ngay dưới khinh khí cầu bằng $2\sqrt{13}$ km.}{d}
  \y{Tại gốc tọa độ $O$ có một radar. Hai máy bay bay theo hai hướng có vectơ chỉ phương $\vec{a}=(-2;1;-3)$ và $\vec{b}=(3;2;1)$. Góc giữa hai hướng bay bằng $120^\circ$.}{d}}

\cauDS[id=C2-B8-D11-TH-064,muc=TH,loigiai={Kiểm tra từng mệnh đề bằng công thức tọa độ. a) Đúng. b) Đúng. c) Sai. d) Đúng.}]{Xét tính đúng sai của các mệnh đề sau:}{%
  \y{Trong không gian $Oxyz$ (đơn vị: km), một máy bay bay thẳng từ $A(1;-3;2)$ đến $B(3;-4;4)$ với tốc độ không đổi $6$ km/h. Thời gian bay từ $A$ đến $B$ là $\dfrac{1}{2}$ giờ.}{d}
  \y{Trong không gian $Oxyz$ (đơn vị: km), một máy bay bay thẳng từ $A(-2;1;-1)$ đến $B(-1;-1;1)$ với tốc độ không đổi $3$ km/h. Thời gian bay từ $A$ đến $B$ là $1$ giờ.}{d}
  \y{Trong không gian $Oxyz$ (đơn vị: km), mặt đất là mặt phẳng $(Oxy)$ và trạm điều khiển đặt tại gốc $O$. Một khinh khí cầu ở vị trí $M(2;4;7)$. Khoảng cách trên mặt đất từ trạm đến điểm ngay dưới khinh khí cầu bằng $\sqrt{19}$ km.}{s}
  \y{Trong không gian $Oxyz$ (đơn vị: m), một tuyến cáp treo thẳng nối hai trạm tại $A(-3;0;-1)$ và $B(-7;4;1)$. Chiều dài tuyến cáp treo bằng $6$ m.}{d}}

\cauTLN[id=C2-B8-D11-TH-065,muc=TH,loigiai={Hình chiếu của $M$ lên mặt đất là $H(6;-6;0)$ nên khoảng cách theo mặt đất từ trạm đến điểm ngay dưới khinh khí cầu là $OH=\sqrt{6^2-6^2}=6\sqrt{2}$ (km).}]{Trong không gian $Oxyz$ (đơn vị: km), mặt đất là mặt phẳng $(Oxy)$ và trạm điều khiển đặt tại gốc $O$. Một khinh khí cầu ở vị trí $M(6;-6;2)$. Tính bình phương khoảng cách trên mặt đất từ trạm đến điểm ngay dưới khinh khí cầu.}{72}

\cauTLN[id=C2-B8-D11-TH-066,muc=TH,loigiai={Hình chiếu của $M$ lên mặt đất là $H(-3;-3;0)$ nên khoảng cách theo mặt đất từ trạm đến điểm ngay dưới khinh khí cầu là $OH=\sqrt{-3^2-3^2}=3\sqrt{2}$ (km).}]{Trong không gian $Oxyz$ (đơn vị: km), mặt đất là mặt phẳng $(Oxy)$ và trạm điều khiển đặt tại gốc $O$. Một khinh khí cầu ở vị trí $M(-3;-3;1)$. Tính bình phương khoảng cách trên mặt đất từ trạm đến điểm ngay dưới khinh khí cầu.}{18}

\cauTLN[id=C2-B8-D11-TH-067,muc=TH,loigiai={Hình chiếu của $M$ lên mặt đất là $H(-2;-6;0)$ nên khoảng cách theo mặt đất từ trạm đến điểm ngay dưới khinh khí cầu là $OH=\sqrt{-2^2-6^2}=2\sqrt{10}$ (km).}]{Trong không gian $Oxyz$ (đơn vị: km), mặt đất là mặt phẳng $(Oxy)$ và trạm điều khiển đặt tại gốc $O$. Một khinh khí cầu ở vị trí $M(-2;-6;6)$. Tính bình phương khoảng cách trên mặt đất từ trạm đến điểm ngay dưới khinh khí cầu.}{40}

\cauTLN[id=C2-B8-D11-TH-068,muc=TH,loigiai={Hình chiếu của $M$ lên mặt đất là $H(0;-3;0)$ nên khoảng cách theo mặt đất từ trạm đến điểm ngay dưới khinh khí cầu là $OH=\sqrt{0^2-3^2}=3$ (km).}]{Trong không gian $Oxyz$ (đơn vị: km), mặt đất là mặt phẳng $(Oxy)$ và trạm điều khiển đặt tại gốc $O$. Một khinh khí cầu ở vị trí $M(0;-3;6)$. Tính bình phương khoảng cách trên mặt đất từ trạm đến điểm ngay dưới khinh khí cầu.}{9}

\cauTLN[id=C2-B8-D11-TH-069,muc=TH,loigiai={Hình chiếu của $M$ lên mặt đất là $H(-1;0;0)$ nên khoảng cách theo mặt đất từ trạm đến điểm ngay dưới khinh khí cầu là $OH=\sqrt{-1^2+0^2}=1$ (km).}]{Trong không gian $Oxyz$ (đơn vị: km), mặt đất là mặt phẳng $(Oxy)$ và trạm điều khiển đặt tại gốc $O$. Một khinh khí cầu ở vị trí $M(-1;0;1)$. Tính bình phương khoảng cách trên mặt đất từ trạm đến điểm ngay dưới khinh khí cầu.}{1}

\cauTLN[id=C2-B8-D11-TH-070,muc=TH,loigiai={$AB=\sqrt{(-4)^2+(4)^2+(2)^2}=6$ (m).}]{Trong không gian $Oxyz$ (đơn vị: m), một tuyến cáp treo thẳng nối hai trạm tại $A(1;1;3)$ và $B(-3;5;5)$. Tính chiều dài tuyến cáp treo (đơn vị: m).}{6}

\cauTLN[id=C2-B8-D11-TH-071,muc=TH,loigiai={$AB=\sqrt{(2)^2+(2)^2+(-1)^2}=3$ (m).}]{Trong không gian $Oxyz$ (đơn vị: m), một tuyến cáp treo thẳng nối hai trạm tại $A(0;0;3)$ và $B(2;2;2)$. Tính chiều dài tuyến cáp treo (đơn vị: m).}{3}

\cauTLN[id=C2-B8-D11-TH-072,muc=TH,loigiai={$AB=\sqrt{(5)^2+(0)^2+(0)^2}=5$ (m).}]{Trong không gian $Oxyz$ (đơn vị: m), một tuyến cáp treo thẳng nối hai trạm tại $A(1;-2;1)$ và $B(6;-2;1)$. Tính chiều dài tuyến cáp treo (đơn vị: m).}{5}

\cauTLN[id=C2-B8-D11-TH-073,muc=TH,loigiai={$AB=\sqrt{(2)^2+(4)^2+(4)^2}=6$ (m).}]{Trong không gian $Oxyz$ (đơn vị: m), một tuyến cáp treo thẳng nối hai trạm tại $A(-2;2;-2)$ và $B(0;6;2)$. Tính chiều dài tuyến cáp treo (đơn vị: m).}{6}

\cauTLN[id=C2-B8-D11-TH-074,muc=TH,loigiai={$AB=\sqrt{(-4)^2+(4)^2+(2)^2}=6$ (m).}]{Trong không gian $Oxyz$ (đơn vị: m), một tuyến cáp treo thẳng nối hai trạm tại $A(0;-1;2)$ và $B(-4;3;4)$. Tính chiều dài tuyến cáp treo (đơn vị: m).}{6}

\cauTLN[id=C2-B8-D11-TH-075,muc=TH,loigiai={$AB=15$ km nên thời gian bay là $\dfrac{15}{3}=5$ giờ.}]{Trong không gian $Oxyz$ (đơn vị: km), một máy bay bay thẳng từ $A(2;0;-1)$ đến $B(14;9;-1)$ với tốc độ không đổi $3$ km/h. Tính thời gian bay từ $A$ đến $B$ (đơn vị: giờ).}{5}

\cauTLN[id=C2-B8-D11-TH-076,muc=TH,loigiai={$AB=12$ km nên thời gian bay là $\dfrac{12}{12}=1$ giờ.}]{Trong không gian $Oxyz$ (đơn vị: km), một máy bay bay thẳng từ $A(-1;3;-1)$ đến $B(3;-5;7)$ với tốc độ không đổi $12$ km/h. Tính thời gian bay từ $A$ đến $B$ (đơn vị: giờ).}{1}

\cauTLN[id=C2-B8-D11-TH-077,muc=TH,loigiai={$AB=6$ km nên thời gian bay là $\dfrac{6}{3}=2$ giờ.}]{Trong không gian $Oxyz$ (đơn vị: km), một máy bay bay thẳng từ $A(-1;-1;-3)$ đến $B(3;3;-5)$ với tốc độ không đổi $3$ km/h. Tính thời gian bay từ $A$ đến $B$ (đơn vị: giờ).}{2}

\cauTLN[id=C2-B8-D11-TH-078,muc=TH,loigiai={$AB=7$ km nên thời gian bay là $\dfrac{7}{7}=1$ giờ.}]{Trong không gian $Oxyz$ (đơn vị: km), một máy bay bay thẳng từ $A(-2;1;1)$ đến $B(4;4;3)$ với tốc độ không đổi $7$ km/h. Tính thời gian bay từ $A$ đến $B$ (đơn vị: giờ).}{1}

\cauTLN[id=C2-B8-D11-TH-079,muc=TH,loigiai={$AB=12$ km nên thời gian bay là $\dfrac{12}{3}=4$ giờ.}]{Trong không gian $Oxyz$ (đơn vị: km), một máy bay bay thẳng từ $A(0;-3;-2)$ đến $B(8;1;6)$ với tốc độ không đổi $3$ km/h. Tính thời gian bay từ $A$ đến $B$ (đơn vị: giờ).}{4}

\cauTLN[id=C2-B8-D11-TH-080,muc=TH,loigiai={$AB=12$ km nên thời gian bay là $\dfrac{12}{3}=4$ giờ.}]{Trong không gian $Oxyz$ (đơn vị: km), một máy bay bay thẳng từ $A(3;-3;3)$ đến $B(11;5;7)$ với tốc độ không đổi $3$ km/h. Tính thời gian bay từ $A$ đến $B$ (đơn vị: giờ).}{4}

\cauTN[id=C2-B8-D11-VD-001,muc=VD,loigiai={Sau $3$ giờ, hai máy bay ở $P_1=(6;0;6)$ và $P_2=(6;3;-3)$ nên khoảng cách là $P_1P_2=\sqrt{90}=3\sqrt{10}$ (km).}]{Trong không gian $Oxyz$ (đơn vị: km), hai máy bay cùng xuất phát lúc $0$ giờ từ $A(3;-3;0)$ và $B(3;0;0)$, bay thẳng với vận tốc không đổi lần lượt là $\vec{u}=(1;1;2)$ và $\vec{v}=(1;1;-1)$ (km/h). Sau $3$ giờ, khoảng cách giữa hai máy bay bằng}
  {$3\sqrt{10}$.}
  {$\sqrt{91}$.}{$\sqrt{89}$.}{$\sqrt{94}$.}

\cauTN[id=C2-B8-D11-VD-002,muc=VD,loigiai={Sau $3$ giờ, hai máy bay ở $P_1=(-2;10;5)$ và $P_2=(-3;4;-1)$ nên khoảng cách là $P_1P_2=\sqrt{73}=\sqrt{73}$ (km).}]{Trong không gian $Oxyz$ (đơn vị: km), hai máy bay cùng xuất phát lúc $0$ giờ từ $A(4;4;-1)$ và $B(-3;-2;-1)$, bay thẳng với vận tốc không đổi lần lượt là $\vec{u}=(-2;2;2)$ và $\vec{v}=(0;2;0)$ (km/h). Sau $3$ giờ, khoảng cách giữa hai máy bay bằng}
  {$\sqrt{73}$.}
  {$\sqrt{74}$.}{$6\sqrt{2}$.}{$\sqrt{77}$.}

\cauTN[id=C2-B8-D11-VD-003,muc=VD,loigiai={Sau $3$ giờ, hai máy bay ở $P_1=(-7;-3;4)$ và $P_2=(5;-4;-2)$ nên khoảng cách là $P_1P_2=\sqrt{181}=\sqrt{181}$ (km).}]{Trong không gian $Oxyz$ (đơn vị: km), hai máy bay cùng xuất phát lúc $0$ giờ từ $A(-4;0;4)$ và $B(2;-1;-2)$, bay thẳng với vận tốc không đổi lần lượt là $\vec{u}=(-1;-1;0)$ và $\vec{v}=(1;-1;0)$ (km/h). Sau $3$ giờ, khoảng cách giữa hai máy bay bằng}
  {$\sqrt{181}$.}
  {$\sqrt{182}$.}{$6\sqrt{5}$.}{$\sqrt{185}$.}

\cauTN[id=C2-B8-D11-VD-004,muc=VD,loigiai={Sau $3$ giờ, hai máy bay ở $P_1=(6;7;-7)$ và $P_2=(6;-5;2)$ nên khoảng cách là $P_1P_2=\sqrt{225}=15$ (km).}]{Trong không gian $Oxyz$ (đơn vị: km), hai máy bay cùng xuất phát lúc $0$ giờ từ $A(3;1;-1)$ và $B(3;1;2)$, bay thẳng với vận tốc không đổi lần lượt là $\vec{u}=(1;2;-2)$ và $\vec{v}=(1;-2;0)$ (km/h). Sau $3$ giờ, khoảng cách giữa hai máy bay bằng}
  {$15$.}
  {$\sqrt{226}$.}{$4\sqrt{14}$.}{$\sqrt{229}$.}

\cauTN[id=C2-B8-D11-VD-005,muc=VD,loigiai={Sau $2$ giờ, hai máy bay ở $P_1=(-5;1;8)$ và $P_2=(-3;0;2)$ nên khoảng cách là $P_1P_2=\sqrt{41}=\sqrt{41}$ (km).}]{Trong không gian $Oxyz$ (đơn vị: km), hai máy bay cùng xuất phát lúc $0$ giờ từ $A(-3;-3;4)$ và $B(-3;0;4)$, bay thẳng với vận tốc không đổi lần lượt là $\vec{u}=(-1;2;2)$ và $\vec{v}=(0;0;-1)$ (km/h). Sau $2$ giờ, khoảng cách giữa hai máy bay bằng}
  {$\sqrt{41}$.}
  {$\sqrt{42}$.}{$2\sqrt{10}$.}{$3\sqrt{5}$.}

\cauTN[id=C2-B8-D11-VD-006,muc=VD,loigiai={Sau $2$ giờ, hai máy bay ở $P_1=(1;-4;0)$ và $P_2=(6;1;-1)$ nên khoảng cách là $P_1P_2=\sqrt{51}=\sqrt{51}$ (km).}]{Trong không gian $Oxyz$ (đơn vị: km), hai máy bay cùng xuất phát lúc $0$ giờ từ $A(-3;0;2)$ và $B(4;-3;-1)$, bay thẳng với vận tốc không đổi lần lượt là $\vec{u}=(2;-2;-1)$ và $\vec{v}=(1;2;0)$ (km/h). Sau $2$ giờ, khoảng cách giữa hai máy bay bằng}
  {$\sqrt{51}$.}
  {$2\sqrt{13}$.}{$5\sqrt{2}$.}{$\sqrt{55}$.}

\cauTN[id=C2-B8-D11-VD-007,muc=VD,loigiai={Sau $3$ giờ, hai máy bay ở $P_1=(10;-3;2)$ và $P_2=(-5;3;0)$ nên khoảng cách là $P_1P_2=\sqrt{265}=\sqrt{265}$ (km).}]{Trong không gian $Oxyz$ (đơn vị: km), hai máy bay cùng xuất phát lúc $0$ giờ từ $A(4;0;-4)$ và $B(-2;0;0)$, bay thẳng với vận tốc không đổi lần lượt là $\vec{u}=(2;-1;2)$ và $\vec{v}=(-1;1;0)$ (km/h). Sau $3$ giờ, khoảng cách giữa hai máy bay bằng}
  {$\sqrt{265}$.}
  {$\sqrt{266}$.}{$2\sqrt{66}$.}{$\sqrt{269}$.}

\cauTN[id=C2-B8-D11-VD-008,muc=VD,loigiai={Sau $3$ giờ, hai máy bay ở $P_1=(-4;-2;-3)$ và $P_2=(7;-2;-3)$ nên khoảng cách là $P_1P_2=\sqrt{121}=11$ (km).}]{Trong không gian $Oxyz$ (đơn vị: km), hai máy bay cùng xuất phát lúc $0$ giờ từ $A(2;-2;0)$ và $B(4;1;-3)$, bay thẳng với vận tốc không đổi lần lượt là $\vec{u}=(-2;0;-1)$ và $\vec{v}=(1;-1;0)$ (km/h). Sau $3$ giờ, khoảng cách giữa hai máy bay bằng}
  {$11$.}
  {$\sqrt{122}$.}{$2\sqrt{30}$.}{$5\sqrt{5}$.}

\cauTN[id=C2-B8-D11-VD-009,muc=VD,loigiai={Sau $3$ giờ, hai máy bay ở $P_1=(-2;4;0)$ và $P_2=(7;-4;1)$ nên khoảng cách là $P_1P_2=\sqrt{146}=\sqrt{146}$ (km).}]{Trong không gian $Oxyz$ (đơn vị: km), hai máy bay cùng xuất phát lúc $0$ giờ từ $A(4;-2;3)$ và $B(4;-4;-2)$, bay thẳng với vận tốc không đổi lần lượt là $\vec{u}=(-2;2;-1)$ và $\vec{v}=(1;0;1)$ (km/h). Sau $3$ giờ, khoảng cách giữa hai máy bay bằng}
  {$\sqrt{146}$.}
  {$7\sqrt{3}$.}{$\sqrt{145}$.}{$5\sqrt{6}$.}

\cauTN[id=C2-B8-D11-VD-010,muc=VD,loigiai={Sau $1$ giờ, hai máy bay ở $P_1=(3;-2;3)$ và $P_2=(-3;-3;-3)$ nên khoảng cách là $P_1P_2=\sqrt{73}=\sqrt{73}$ (km).}]{Trong không gian $Oxyz$ (đơn vị: km), hai máy bay cùng xuất phát lúc $0$ giờ từ $A(3;-4;4)$ và $B(-1;-4;-3)$, bay thẳng với vận tốc không đổi lần lượt là $\vec{u}=(0;2;-1)$ và $\vec{v}=(-2;1;0)$ (km/h). Sau $1$ giờ, khoảng cách giữa hai máy bay bằng}
  {$\sqrt{73}$.}
  {$\sqrt{74}$.}{$6\sqrt{2}$.}{$\sqrt{77}$.}

\cauTN[id=C2-B8-D11-VD-011,muc=VD,loigiai={Sau $1$ giờ, hai máy bay ở $P_1=(-1;6;0)$ và $P_2=(-2;2;-1)$ nên khoảng cách là $P_1P_2=\sqrt{18}=3\sqrt{2}$ (km).}]{Trong không gian $Oxyz$ (đơn vị: km), hai máy bay cùng xuất phát lúc $0$ giờ từ $A(-2;4;-1)$ và $B(-3;4;-3)$, bay thẳng với vận tốc không đổi lần lượt là $\vec{u}=(1;2;1)$ và $\vec{v}=(1;-2;2)$ (km/h). Sau $1$ giờ, khoảng cách giữa hai máy bay bằng}
  {$3\sqrt{2}$.}
  {$\sqrt{19}$.}{$\sqrt{17}$.}{$\sqrt{22}$.}

\cauTN[id=C2-B8-D11-VD-012,muc=VD,loigiai={Sau $3$ giờ, hai máy bay ở $P_1=(-3;-10;-3)$ và $P_2=(-4;6;3)$ nên khoảng cách là $P_1P_2=\sqrt{293}=\sqrt{293}$ (km).}]{Trong không gian $Oxyz$ (đơn vị: km), hai máy bay cùng xuất phát lúc $0$ giờ từ $A(-3;-4;3)$ và $B(-1;3;3)$, bay thẳng với vận tốc không đổi lần lượt là $\vec{u}=(0;-2;-2)$ và $\vec{v}=(-1;1;0)$ (km/h). Sau $3$ giờ, khoảng cách giữa hai máy bay bằng}
  {$\sqrt{293}$.}
  {$7\sqrt{6}$.}{$2\sqrt{73}$.}{$3\sqrt{33}$.}

\cauTN[id=C2-B8-D11-VD-013,muc=VD,loigiai={$AB=1$, $BC=7$ nên quãng đường bay là $1+7=8$ (km).}]{Trong không gian $Oxyz$ (đơn vị: km), một drone bay từ $A(2;2;1)$ đến $B(3;2;1)$ rồi bay thẳng đến $C(5;8;4)$. Tổng quãng đường drone đã bay (đơn vị: km) là}
  {$8$.}
  {$9$.}{$7$.}{$10$.}

\cauTN[id=C2-B8-D11-VD-014,muc=VD,loigiai={$AB=3$, $BC=5$ nên quãng đường bay là $3+5=8$ (km).}]{Trong không gian $Oxyz$ (đơn vị: km), một drone bay từ $A(-3;3;-1)$ đến $B(-1;4;1)$ rồi bay thẳng đến $C(3;4;4)$. Tổng quãng đường drone đã bay (đơn vị: km) là}
  {$8$.}
  {$9$.}{$7$.}{$10$.}

\cauTN[id=C2-B8-D11-VD-015,muc=VD,loigiai={$AB=7$, $BC=3$ nên quãng đường bay là $7+3=10$ (km).}]{Trong không gian $Oxyz$ (đơn vị: km), một drone bay từ $A(-2;2;2)$ đến $B(4;5;4)$ rồi bay thẳng đến $C(6;6;6)$. Tổng quãng đường drone đã bay (đơn vị: km) là}
  {$10$.}
  {$11$.}{$9$.}{$12$.}

\cauTN[id=C2-B8-D11-VD-016,muc=VD,loigiai={$AB=5$, $BC=5$ nên quãng đường bay là $5+5=10$ (km).}]{Trong không gian $Oxyz$ (đơn vị: km), một drone bay từ $A(0;3;2)$ đến $B(3;7;2)$ rồi bay thẳng đến $C(3;7;7)$. Tổng quãng đường drone đã bay (đơn vị: km) là}
  {$10$.}
  {$11$.}{$9$.}{$7$.}

\cauTN[id=C2-B8-D11-VD-017,muc=VD,loigiai={$AB=5$, $BC=5$ nên quãng đường bay là $5+5=10$ (km).}]{Trong không gian $Oxyz$ (đơn vị: km), một drone bay từ $A(3;3;0)$ đến $B(6;7;0)$ rồi bay thẳng đến $C(10;10;0)$. Tổng quãng đường drone đã bay (đơn vị: km) là}
  {$10$.}
  {$11$.}{$9$.}{$12$.}

\cauTN[id=C2-B8-D11-VD-018,muc=VD,loigiai={$AB=3$, $BC=3$ nên quãng đường bay là $3+3=6$ (km).}]{Trong không gian $Oxyz$ (đơn vị: km), một drone bay từ $A(-3;-1;-3)$ đến $B(-5;1;-2)$ rồi bay thẳng đến $C(-3;0;0)$. Tổng quãng đường drone đã bay (đơn vị: km) là}
  {$6$.}
  {$7$.}{$5$.}{$3$.}

\cauTN[id=C2-B8-D11-VD-019,muc=VD,loigiai={$AB=11$, $BC=7$ nên quãng đường bay là $11+7=18$ (km).}]{Trong không gian $Oxyz$ (đơn vị: km), một drone bay từ $A(-3;-3;-2)$ đến $B(3;3;5)$ rồi bay thẳng đến $C(6;5;11)$. Tổng quãng đường drone đã bay (đơn vị: km) là}
  {$18$.}
  {$19$.}{$17$.}{$20$.}

\cauTN[id=C2-B8-D11-VD-020,muc=VD,loigiai={$AB=5$, $BC=5$ nên quãng đường bay là $5+5=10$ (km).}]{Trong không gian $Oxyz$ (đơn vị: km), một drone bay từ $A(1;-1;2)$ đến $B(1;-1;7)$ rồi bay thẳng đến $C(4;3;7)$. Tổng quãng đường drone đã bay (đơn vị: km) là}
  {$10$.}
  {$11$.}{$9$.}{$7$.}

\cauTN[id=C2-B8-D11-VD-021,muc=VD,loigiai={$AB=3$, $BC=1$ nên quãng đường bay là $3+1=4$ (km).}]{Trong không gian $Oxyz$ (đơn vị: km), một drone bay từ $A(-1;-3;0)$ đến $B(1;-4;2)$ rồi bay thẳng đến $C(1;-3;2)$. Tổng quãng đường drone đã bay (đơn vị: km) là}
  {$4$.}
  {$5$.}{$3$.}{$6$.}

\cauTN[id=C2-B8-D11-VD-022,muc=VD,loigiai={$AB=1$, $BC=5$ nên quãng đường bay là $1+5=6$ (km).}]{Trong không gian $Oxyz$ (đơn vị: km), một drone bay từ $A(3;-3;2)$ đến $B(3;-3;3)$ rồi bay thẳng đến $C(6;-3;7)$. Tổng quãng đường drone đã bay (đơn vị: km) là}
  {$6$.}
  {$7$.}{$5$.}{$8$.}

\cauTN[id=C2-B8-D11-VD-023,muc=VD,loigiai={$AB=5$, $BC=7$ nên quãng đường bay là $5+7=12$ (km).}]{Trong không gian $Oxyz$ (đơn vị: km), một drone bay từ $A(1;2;1)$ đến $B(1;5;5)$ rồi bay thẳng đến $C(3;11;8)$. Tổng quãng đường drone đã bay (đơn vị: km) là}
  {$12$.}
  {$13$.}{$11$.}{$14$.}

\cauTN[id=C2-B8-D11-VD-024,muc=VD,loigiai={$AB=5$, $BC=5$ nên quãng đường bay là $5+5=10$ (km).}]{Trong không gian $Oxyz$ (đơn vị: km), một drone bay từ $A(2;3;3)$ đến $B(6;3;6)$ rồi bay thẳng đến $C(6;7;9)$. Tổng quãng đường drone đã bay (đơn vị: km) là}
  {$10$.}
  {$11$.}{$9$.}{$8$.}

\cauTN[id=C2-B8-D11-VD-025,muc=VD,loigiai={Vì $AM:MB=2:1$ nên $\overrightarrow{AM}=\dfrac{2}{3}\overrightarrow{AB}$, suy ra $M=(-2;-2;-2)$.}]{Trong không gian $Oxyz$ (đơn vị: km), hai trạm ở $A(4;2;-2)$ và $B(-5;-4;-2)$. Một trạm chuyển tiếp $M$ nằm trên đoạn $AB$ sao cho $AM:MB=2:1$. Tọa độ của trạm chuyển tiếp $M$ là}
  {$(-2;-2;-2)$.}
  {$(1;0;-2)$.}{$(-2;-2;-1)$.}{$(-2;-2;-4)$.}

\cauTN[id=C2-B8-D11-VD-026,muc=VD,loigiai={Vì $AM:MB=3:2$ nên $\overrightarrow{AM}=\dfrac{3}{5}\overrightarrow{AB}$, suy ra $M=(-2;-3;2)$.}]{Trong không gian $Oxyz$ (đơn vị: km), hai trạm ở $A(-5;0;5)$ và $B(0;-5;0)$. Một trạm chuyển tiếp $M$ nằm trên đoạn $AB$ sao cho $AM:MB=3:2$. Tọa độ của trạm chuyển tiếp $M$ là}
  {$(-2;-3;2)$.}
  {$(2;-3;-2)$.}{$(-2;-2;2)$.}{$(-2;-1;2)$.}

\cauTN[id=C2-B8-D11-VD-027,muc=VD,loigiai={Vì $AM:MB=3:2$ nên $\overrightarrow{AM}=\dfrac{3}{5}\overrightarrow{AB}$, suy ra $M=(-1;0;3)$.}]{Trong không gian $Oxyz$ (đơn vị: km), hai trạm ở $A(-4;0;3)$ và $B(1;0;3)$. Một trạm chuyển tiếp $M$ nằm trên đoạn $AB$ sao cho $AM:MB=3:2$. Tọa độ của trạm chuyển tiếp $M$ là}
  {$(-1;0;3)$.}
  {$(1;0;-3)$.}{$(0;0;3)$.}{$(-1;0;5)$.}

\cauTN[id=C2-B8-D11-VD-028,muc=VD,loigiai={Vì $AM:MB=2:1$ nên $\overrightarrow{AM}=\dfrac{2}{3}\overrightarrow{AB}$, suy ra $M=(3;2;3)$.}]{Trong không gian $Oxyz$ (đơn vị: km), hai trạm ở $A(3;-2;-1)$ và $B(3;4;5)$. Một trạm chuyển tiếp $M$ nằm trên đoạn $AB$ sao cho $AM:MB=2:1$. Tọa độ của trạm chuyển tiếp $M$ là}
  {$(3;2;3)$.}
  {$(3;4;3)$.}{$(3;0;1)$.}{$(5;2;3)$.}

\cauTN[id=C2-B8-D11-VD-029,muc=VD,loigiai={Vì $AM:MB=1:2$ nên $\overrightarrow{AM}=\dfrac{1}{3}\overrightarrow{AB}$, suy ra $M=(3;-2;-1)$.}]{Trong không gian $Oxyz$ (đơn vị: km), hai trạm ở $A(3;-4;0)$ và $B(3;2;-3)$. Một trạm chuyển tiếp $M$ nằm trên đoạn $AB$ sao cho $AM:MB=1:2$. Tọa độ của trạm chuyển tiếp $M$ là}
  {$(3;-2;-1)$.}
  {$(2;-2;-1)$.}{$(3;-3;-1)$.}{$(3;2;-1)$.}

\cauTN[id=C2-B8-D11-VD-030,muc=VD,loigiai={Vì $AM:MB=2:1$ nên $\overrightarrow{AM}=\dfrac{2}{3}\overrightarrow{AB}$, suy ra $M=(-3;-1;1)$.}]{Trong không gian $Oxyz$ (đơn vị: km), hai trạm ở $A(-3;-3;1)$ và $B(-3;0;1)$. Một trạm chuyển tiếp $M$ nằm trên đoạn $AB$ sao cho $AM:MB=2:1$. Tọa độ của trạm chuyển tiếp $M$ là}
  {$(-3;-1;1)$.}
  {$(1;-1;-3)$.}{$(-3;-1;2)$.}{$(-3;1;1)$.}

\cauTN[id=C2-B8-D11-VD-031,muc=VD,loigiai={Vì $AM:MB=1:2$ nên $\overrightarrow{AM}=\dfrac{1}{3}\overrightarrow{AB}$, suy ra $M=(1;-1;0)$.}]{Trong không gian $Oxyz$ (đơn vị: km), hai trạm ở $A(0;1;2)$ và $B(3;-5;-4)$. Một trạm chuyển tiếp $M$ nằm trên đoạn $AB$ sao cho $AM:MB=1:2$. Tọa độ của trạm chuyển tiếp $M$ là}
  {$(1;-1;0)$.}
  {$(1;1;0)$.}{$(1;-1;1)$.}{$(1;0;-1)$.}

\cauTN[id=C2-B8-D11-VD-032,muc=VD,loigiai={Vì $AM:MB=3:1$ nên $\overrightarrow{AM}=\dfrac{3}{4}\overrightarrow{AB}$, suy ra $M=(2;3;-1)$.}]{Trong không gian $Oxyz$ (đơn vị: km), hai trạm ở $A(2;3;5)$ và $B(2;3;-3)$. Một trạm chuyển tiếp $M$ nằm trên đoạn $AB$ sao cho $AM:MB=3:1$. Tọa độ của trạm chuyển tiếp $M$ là}
  {$(2;3;-1)$.}
  {$(3;3;-1)$.}{$(2;-1;3)$.}{$(1;3;-1)$.}

\cauTN[id=C2-B8-D11-VD-033,muc=VD,loigiai={Vì $AM:MB=3:1$ nên $\overrightarrow{AM}=\dfrac{3}{4}\overrightarrow{AB}$, suy ra $M=(-1;-3;1)$.}]{Trong không gian $Oxyz$ (đơn vị: km), hai trạm ở $A(-1;0;-2)$ và $B(-1;-4;2)$. Một trạm chuyển tiếp $M$ nằm trên đoạn $AB$ sao cho $AM:MB=3:1$. Tọa độ của trạm chuyển tiếp $M$ là}
  {$(-1;-3;1)$.}
  {$(-1;-1;1)$.}{$(-1;-3;0)$.}{$(-1;-3;-1)$.}

\cauTN[id=C2-B8-D11-VD-034,muc=VD,loigiai={Vì $AM:MB=2:1$ nên $\overrightarrow{AM}=\dfrac{2}{3}\overrightarrow{AB}$, suy ra $M=(-3;0;-2)$.}]{Trong không gian $Oxyz$ (đơn vị: km), hai trạm ở $A(1;0;-4)$ và $B(-5;0;-1)$. Một trạm chuyển tiếp $M$ nằm trên đoạn $AB$ sao cho $AM:MB=2:1$. Tọa độ của trạm chuyển tiếp $M$ là}
  {$(-3;0;-2)$.}
  {$(-3;0;-1)$.}{$(-2;0;-2)$.}{$(0;-3;-2)$.}

\cauTN[id=C2-B8-D11-VD-035,muc=VD,loigiai={Vì $AM:MB=1:3$ nên $\overrightarrow{AM}=\dfrac{1}{4}\overrightarrow{AB}$, suy ra $M=(0;-1;4)$.}]{Trong không gian $Oxyz$ (đơn vị: km), hai trạm ở $A(0;-3;4)$ và $B(0;5;4)$. Một trạm chuyển tiếp $M$ nằm trên đoạn $AB$ sao cho $AM:MB=1:3$. Tọa độ của trạm chuyển tiếp $M$ là}
  {$(0;-1;4)$.}
  {$(0;-3;4)$.}{$(0;1;-4)$.}{$(1;-1;4)$.}

\cauTN[id=C2-B8-D11-VD-036,muc=VD,loigiai={Vì $AM:MB=1:2$ nên $\overrightarrow{AM}=\dfrac{1}{3}\overrightarrow{AB}$, suy ra $M=(2;1;1)$.}]{Trong không gian $Oxyz$ (đơn vị: km), hai trạm ở $A(2;3;-1)$ và $B(2;-3;5)$. Một trạm chuyển tiếp $M$ nằm trên đoạn $AB$ sao cho $AM:MB=1:2$. Tọa độ của trạm chuyển tiếp $M$ là}
  {$(2;1;1)$.}
  {$(1;2;1)$.}{$(3;1;1)$.}{$(2;1;2)$.}

\cauDS[id=C2-B8-D11-VD-037,muc=VD,loigiai={Kiểm tra từng mệnh đề bằng công thức tọa độ. a) Sai. b) Đúng. c) Sai. d) Đúng.}]{Xét tính đúng sai của các mệnh đề sau:}{%
  \y{Trong không gian $Oxyz$ (đơn vị: km), một drone bay từ $A(0;0;-2)$ đến $B(3;4;-2)$ rồi bay thẳng đến $C(4;6;0)$. Tổng quãng đường drone bay là $10$ km.}{s}
  \y{Trong không gian $Oxyz$ (đơn vị: km), hai máy bay cùng xuất phát lúc $0$ giờ từ $A(4;0;-4)$ và $B(4;2;-1)$, bay thẳng với vận tốc không đổi lần lượt là $\vec{u}=(-2;-2;1)$ và $\vec{v}=(0;-2;2)$ (km/h). Sau $2$ giờ, khoảng cách giữa hai máy bay bằng $3\sqrt{5}$ km.}{d}
  \y{Trong không gian $Oxyz$ (đơn vị: km), hai máy bay cùng xuất phát lúc $0$ giờ từ $A(-2;-4;4)$ và $B(-2;1;-4)$, bay thẳng với vận tốc không đổi lần lượt là $\vec{u}=(-1;2;-2)$ và $\vec{v}=(-1;0;0)$ (km/h). Sau $2$ giờ, khoảng cách giữa hai máy bay bằng $\sqrt{21}$ km.}{s}
  \y{Trong không gian $Oxyz$ (đơn vị: km), hai trạm ở $A(-1;2;2)$ và $B(2;2;2)$. Một trạm chuyển tiếp $M$ nằm trên đoạn $AB$ sao cho $AM:MB=1:2$. Trạm chuyển tiếp $M$ có tọa độ $(0;2;2)$.}{d}}

\cauDS[id=C2-B8-D11-VD-038,muc=VD,loigiai={Kiểm tra từng mệnh đề bằng công thức tọa độ. a) Đúng. b) Đúng. c) Sai. d) Đúng.}]{Xét tính đúng sai của các mệnh đề sau:}{%
  \y{Tại gốc tọa độ $O$ có một radar. Hai máy bay bay theo hai hướng có vectơ chỉ phương $\vec{a}=(-3;1;0)$ và $\vec{b}=(-1;-3;-3)$. Góc giữa hai hướng bay bằng $90^\circ$.}{d}
  \y{Trong không gian $Oxyz$ (đơn vị: km), một drone bay từ $A(-1;-2;1)$ đến $B(-1;1;5)$ rồi bay thẳng đến $C(3;4;5)$. Tổng quãng đường drone bay là $10$ km.}{d}
  \y{Trong không gian $Oxyz$ (đơn vị: km), hai máy bay cùng xuất phát lúc $0$ giờ từ $A(-4;3;1)$ và $B(3;-3;1)$, bay thẳng với vận tốc không đổi lần lượt là $\vec{u}=(1;1;0)$ và $\vec{v}=(2;1;1)$ (km/h). Sau $2$ giờ, khoảng cách giữa hai máy bay bằng $\sqrt{122}$ km.}{s}
  \y{Trong không gian $Oxyz$ (đơn vị: km), một drone bay từ $A(0;-2;2)$ đến $B(3;-2;6)$ rồi bay thẳng đến $C(3;1;10)$. Tổng quãng đường drone bay là $10$ km.}{d}}

\cauDS[id=C2-B8-D11-VD-039,muc=VD,loigiai={Kiểm tra từng mệnh đề bằng công thức tọa độ. a) Đúng. b) Đúng. c) Sai. d) Đúng.}]{Xét tính đúng sai của các mệnh đề sau:}{%
  \y{Tại gốc tọa độ $O$ có một radar. Hai máy bay bay theo hai hướng có vectơ chỉ phương $\vec{a}=(0;-3;-3)$ và $\vec{b}=(-1;1;-1)$. Góc giữa hai hướng bay bằng $90^\circ$.}{d}
  \y{Trong không gian $Oxyz$ (đơn vị: km), hai trạm ở $A(5;4;-3)$ và $B(-4;-2;0)$. Một trạm chuyển tiếp $M$ nằm trên đoạn $AB$ sao cho $AM:MB=1:2$. Trạm chuyển tiếp $M$ có tọa độ $(2;2;-2)$.}{d}
  \y{Trong không gian $Oxyz$ (đơn vị: km), hai máy bay cùng xuất phát lúc $0$ giờ từ $A(-2;4;3)$ và $B(-4;3;-2)$, bay thẳng với vận tốc không đổi lần lượt là $\vec{u}=(0;-2;2)$ và $\vec{v}=(-1;-2;1)$ (km/h). Sau $2$ giờ, khoảng cách giữa hai máy bay bằng $\sqrt{67}$ km.}{s}
  \y{Trong không gian $Oxyz$ (đơn vị: km), hai máy bay cùng xuất phát lúc $0$ giờ từ $A(-1;3;3)$ và $B(-3;-4;2)$, bay thẳng với vận tốc không đổi lần lượt là $\vec{u}=(-2;1;2)$ và $\vec{v}=(0;-1;2)$ (km/h). Sau $3$ giờ, khoảng cách giữa hai máy bay bằng $\sqrt{186}$ km.}{d}}

\cauDS[id=C2-B8-D11-VD-040,muc=VD,loigiai={Kiểm tra từng mệnh đề bằng công thức tọa độ. a) Sai. b) Sai. c) Đúng. d) Sai.}]{Xét tính đúng sai của các mệnh đề sau:}{%
  \y{Trong không gian $Oxyz$ (đơn vị: km), hai trạm ở $A(4;3;-1)$ và $B(-5;3;-1)$. Một trạm chuyển tiếp $M$ nằm trên đoạn $AB$ sao cho $AM:MB=1:2$. Trạm chuyển tiếp $M$ có tọa độ $(-1;3;-1)$.}{s}
  \y{Tại gốc tọa độ $O$ có một radar. Hai máy bay bay theo hai hướng có vectơ chỉ phương $\vec{a}=(3;0;-3)$ và $\vec{b}=(2;0;0)$. Góc giữa hai hướng bay bằng $30^\circ$.}{s}
  \y{Trong không gian $Oxyz$ (đơn vị: km), một drone bay từ $A(1;2;-1)$ đến $B(1;3;-1)$ rồi bay thẳng đến $C(2;1;1)$. Tổng quãng đường drone bay là $4$ km.}{d}
  \y{Trong không gian $Oxyz$ (đơn vị: km), một drone bay từ $A(-2;-2;1)$ đến $B(-4;0;2)$ rồi bay thẳng đến $C(-4;0;3)$. Tổng quãng đường drone bay là $3$ km.}{s}}

\cauDS[id=C2-B8-D11-VD-041,muc=VD,loigiai={Kiểm tra từng mệnh đề bằng công thức tọa độ. a) Sai. b) Đúng. c) Đúng. d) Sai.}]{Xét tính đúng sai của các mệnh đề sau:}{%
  \y{Trong không gian $Oxyz$ (đơn vị: km), hai trạm ở $A(-5;5;2)$ và $B(4;2;2)$. Một trạm chuyển tiếp $M$ nằm trên đoạn $AB$ sao cho $AM:MB=2:1$. Trạm chuyển tiếp $M$ có tọa độ $(1;-3;2)$.}{s}
  \y{Trong không gian $Oxyz$ (đơn vị: km), hai máy bay cùng xuất phát lúc $0$ giờ từ $A(-1;-4;0)$ và $B(-4;-2;-2)$, bay thẳng với vận tốc không đổi lần lượt là $\vec{u}=(0;1;2)$ và $\vec{v}=(1;0;1)$ (km/h). Sau $2$ giờ, khoảng cách giữa hai máy bay bằng $\sqrt{17}$ km.}{d}
  \y{Trong không gian $Oxyz$ (đơn vị: km), hai trạm ở $A(0;-3;-1)$ và $B(3;-3;-1)$. Một trạm chuyển tiếp $M$ nằm trên đoạn $AB$ sao cho $AM:MB=2:1$. Trạm chuyển tiếp $M$ có tọa độ $(2;-3;-1)$.}{d}
  \y{Trong không gian $Oxyz$ (đơn vị: km), một drone bay từ $A(-3;-1;0)$ đến $B(-1;1;1)$ rồi bay thẳng đến $C(1;3;0)$. Tổng quãng đường drone bay là $5$ km.}{s}}

\cauDS[id=C2-B8-D11-VD-042,muc=VD,loigiai={Kiểm tra từng mệnh đề bằng công thức tọa độ. a) Đúng. b) Sai. c) Sai. d) Đúng.}]{Xét tính đúng sai của các mệnh đề sau:}{%
  \y{Trong không gian $Oxyz$ (đơn vị: km), hai máy bay cùng xuất phát lúc $0$ giờ từ $A(1;2;-2)$ và $B(0;-2;2)$, bay thẳng với vận tốc không đổi lần lượt là $\vec{u}=(-1;1;0)$ và $\vec{v}=(0;-1;-1)$ (km/h). Sau $2$ giờ, khoảng cách giữa hai máy bay bằng $\sqrt{69}$ km.}{d}
  \y{Trong không gian $Oxyz$ (đơn vị: km), hai máy bay cùng xuất phát lúc $0$ giờ từ $A(2;-1;4)$ và $B(3;3;0)$, bay thẳng với vận tốc không đổi lần lượt là $\vec{u}=(-1;2;0)$ và $\vec{v}=(0;1;-2)$ (km/h). Sau $3$ giờ, khoảng cách giữa hai máy bay bằng $11$ km.}{s}
  \y{Trong không gian $Oxyz$ (đơn vị: km), một drone bay từ $A(-3;0;2)$ đến $B(1;4;9)$ rồi bay thẳng đến $C(4;8;9)$. Tổng quãng đường drone bay là $16$ km.}{s}
  \y{Tại gốc tọa độ $O$ có một radar. Hai máy bay bay theo hai hướng có vectơ chỉ phương $\vec{a}=(-3;0;1)$ và $\vec{b}=(-2;0;-1)$. Góc giữa hai hướng bay bằng $45^\circ$.}{d}}

\cauDS[id=C2-B8-D11-VD-043,muc=VD,loigiai={Kiểm tra từng mệnh đề bằng công thức tọa độ. a) Đúng. b) Sai. c) Sai. d) Đúng.}]{Xét tính đúng sai của các mệnh đề sau:}{%
  \y{Tại gốc tọa độ $O$ có một radar. Hai máy bay bay theo hai hướng có vectơ chỉ phương $\vec{a}=(3;-3;0)$ và $\vec{b}=(2;2;0)$. Góc giữa hai hướng bay bằng $90^\circ$.}{d}
  \y{Trong không gian $Oxyz$ (đơn vị: km), một drone bay từ $A(-1;-1;-2)$ đến $B(0;-3;0)$ rồi bay thẳng đến $C(4;1;7)$. Tổng quãng đường drone bay là $13$ km.}{s}
  \y{Tại gốc tọa độ $O$ có một radar. Hai máy bay bay theo hai hướng có vectơ chỉ phương $\vec{a}=(-1;0;-1)$ và $\vec{b}=(3;-3;0)$. Góc giữa hai hướng bay bằng $135^\circ$.}{s}
  \y{Trong không gian $Oxyz$ (đơn vị: km), hai máy bay cùng xuất phát lúc $0$ giờ từ $A(4;0;4)$ và $B(-1;-1;3)$, bay thẳng với vận tốc không đổi lần lượt là $\vec{u}=(0;-2;1)$ và $\vec{v}=(1;1;1)$ (km/h). Sau $2$ giờ, khoảng cách giữa hai máy bay bằng $\sqrt{35}$ km.}{d}}

\cauDS[id=C2-B8-D11-VD-044,muc=VD,loigiai={Kiểm tra từng mệnh đề bằng công thức tọa độ. a) Sai. b) Sai. c) Đúng. d) Sai.}]{Xét tính đúng sai của các mệnh đề sau:}{%
  \y{Trong không gian $Oxyz$ (đơn vị: km), một drone bay từ $A(-1;-3;0)$ đến $B(-1;-3;1)$ rồi bay thẳng đến $C(2;-1;7)$. Tổng quãng đường drone bay là $9$ km.}{s}
  \y{Trong không gian $Oxyz$ (đơn vị: km), hai trạm ở $A(-5;5;1)$ và $B(4;5;-5)$. Một trạm chuyển tiếp $M$ nằm trên đoạn $AB$ sao cho $AM:MB=1:2$. Trạm chuyển tiếp $M$ có tọa độ $(-1;5;-2)$.}{s}
  \y{Tại gốc tọa độ $O$ có một radar. Hai máy bay bay theo hai hướng có vectơ chỉ phương $\vec{a}=(0;1;-2)$ và $\vec{b}=(0;-1;-3)$. Góc giữa hai hướng bay bằng $45^\circ$.}{d}
  \y{Trong không gian $Oxyz$ (đơn vị: km), hai trạm ở $A(0;3;-2)$ và $B(4;3;2)$. Một trạm chuyển tiếp $M$ nằm trên đoạn $AB$ sao cho $AM:MB=3:1$. Trạm chuyển tiếp $M$ có tọa độ $(3;4;1)$.}{s}}

\cauDS[id=C2-B8-D11-VD-045,muc=VD,loigiai={Kiểm tra từng mệnh đề bằng công thức tọa độ. a) Sai. b) Sai. c) Đúng. d) Đúng.}]{Xét tính đúng sai của các mệnh đề sau:}{%
  \y{Tại gốc tọa độ $O$ có một radar. Hai máy bay bay theo hai hướng có vectơ chỉ phương $\vec{a}=(1;1;0)$ và $\vec{b}=(1;-1;0)$. Góc giữa hai hướng bay bằng $60^\circ$.}{s}
  \y{Trong không gian $Oxyz$ (đơn vị: km), hai trạm ở $A(1;-1;-5)$ và $B(-5;-1;-2)$. Một trạm chuyển tiếp $M$ nằm trên đoạn $AB$ sao cho $AM:MB=1:2$. Trạm chuyển tiếp $M$ có tọa độ $(-3;-1;-4)$.}{s}
  \y{Trong không gian $Oxyz$ (đơn vị: km), một drone bay từ $A(-3;0;3)$ đến $B(-2;0;3)$ rồi bay thẳng đến $C(3;0;3)$. Tổng quãng đường drone bay là $6$ km.}{d}
  \y{Trong không gian $Oxyz$ (đơn vị: km), một drone bay từ $A(-3;2;-3)$ đến $B(1;5;-3)$ rồi bay thẳng đến $C(3;6;-1)$. Tổng quãng đường drone bay là $8$ km.}{d}}

\cauDS[id=C2-B8-D11-VD-046,muc=VD,loigiai={Kiểm tra từng mệnh đề bằng công thức tọa độ. a) Đúng. b) Đúng. c) Sai. d) Sai.}]{Xét tính đúng sai của các mệnh đề sau:}{%
  \y{Trong không gian $Oxyz$ (đơn vị: km), hai máy bay cùng xuất phát lúc $0$ giờ từ $A(4;3;-1)$ và $B(4;-2;3)$, bay thẳng với vận tốc không đổi lần lượt là $\vec{u}=(2;1;2)$ và $\vec{v}=(2;-2;0)$ (km/h). Sau $3$ giờ, khoảng cách giữa hai máy bay bằng $10\sqrt{2}$ km.}{d}
  \y{Trong không gian $Oxyz$ (đơn vị: km), một drone bay từ $A(-3;2;2)$ đến $B(1;2;5)$ rồi bay thẳng đến $C(1;2;10)$. Tổng quãng đường drone bay là $10$ km.}{d}
  \y{Trong không gian $Oxyz$ (đơn vị: km), hai máy bay cùng xuất phát lúc $0$ giờ từ $A(0;-1;-1)$ và $B(-4;2;-1)$, bay thẳng với vận tốc không đổi lần lượt là $\vec{u}=(-2;0;0)$ và $\vec{v}=(0;-1;0)$ (km/h). Sau $2$ giờ, khoảng cách giữa hai máy bay bằng $\sqrt{2}$ km.}{s}
  \y{Tại gốc tọa độ $O$ có một radar. Hai máy bay bay theo hai hướng có vectơ chỉ phương $\vec{a}=(0;-2;0)$ và $\vec{b}=(0;2;2)$. Góc giữa hai hướng bay bằng $45^\circ$.}{s}}

\cauDS[id=C2-B8-D11-VD-047,muc=VD,loigiai={Kiểm tra từng mệnh đề bằng công thức tọa độ. a) Sai. b) Sai. c) Đúng. d) Đúng.}]{Xét tính đúng sai của các mệnh đề sau:}{%
  \y{Tại gốc tọa độ $O$ có một radar. Hai máy bay bay theo hai hướng có vectơ chỉ phương $\vec{a}=(-1;1;1)$ và $\vec{b}=(-1;-2;1)$. Góc giữa hai hướng bay bằng $60^\circ$.}{s}
  \y{Trong không gian $Oxyz$ (đơn vị: km), hai máy bay cùng xuất phát lúc $0$ giờ từ $A(-4;-1;-3)$ và $B(0;2;0)$, bay thẳng với vận tốc không đổi lần lượt là $\vec{u}=(2;-2;-1)$ và $\vec{v}=(-2;-2;-2)$ (km/h). Sau $1$ giờ, khoảng cách giữa hai máy bay bằng $\sqrt{17}$ km.}{s}
  \y{Trong không gian $Oxyz$ (đơn vị: km), hai trạm ở $A(5;-2;4)$ và $B(-4;1;1)$. Một trạm chuyển tiếp $M$ nằm trên đoạn $AB$ sao cho $AM:MB=2:1$. Trạm chuyển tiếp $M$ có tọa độ $(-1;0;2)$.}{d}
  \y{Trong không gian $Oxyz$ (đơn vị: km), hai máy bay cùng xuất phát lúc $0$ giờ từ $A(-2;4;1)$ và $B(1;0;4)$, bay thẳng với vận tốc không đổi lần lượt là $\vec{u}=(1;0;2)$ và $\vec{v}=(-1;1;-2)$ (km/h). Sau $2$ giờ, khoảng cách giữa hai máy bay bằng $\sqrt{30}$ km.}{d}}

\cauDS[id=C2-B8-D11-VD-048,muc=VD,loigiai={Kiểm tra từng mệnh đề bằng công thức tọa độ. a) Sai. b) Đúng. c) Sai. d) Sai.}]{Xét tính đúng sai của các mệnh đề sau:}{%
  \y{Trong không gian $Oxyz$ (đơn vị: km), hai trạm ở $A(0;1;0)$ và $B(0;1;5)$. Một trạm chuyển tiếp $M$ nằm trên đoạn $AB$ sao cho $AM:MB=2:3$. Trạm chuyển tiếp $M$ có tọa độ $(1;1;2)$.}{s}
  \y{Trong không gian $Oxyz$ (đơn vị: km), hai máy bay cùng xuất phát lúc $0$ giờ từ $A(2;0;-1)$ và $B(4;-1;-2)$, bay thẳng với vận tốc không đổi lần lượt là $\vec{u}=(1;2;1)$ và $\vec{v}=(-2;1;-2)$ (km/h). Sau $2$ giờ, khoảng cách giữa hai máy bay bằng $\sqrt{74}$ km.}{d}
  \y{Trong không gian $Oxyz$ (đơn vị: km), một drone bay từ $A(-3;2;1)$ đến $B(2;2;1)$ rồi bay thẳng đến $C(2;2;6)$. Tổng quãng đường drone bay là $9$ km.}{s}
  \y{Trong không gian $Oxyz$ (đơn vị: km), hai máy bay cùng xuất phát lúc $0$ giờ từ $A(-1;-1;1)$ và $B(1;-2;1)$, bay thẳng với vận tốc không đổi lần lượt là $\vec{u}=(2;2;2)$ và $\vec{v}=(0;-2;2)$ (km/h). Sau $1$ giờ, khoảng cách giữa hai máy bay bằng $\sqrt{29}$ km.}{s}}

\cauTLN[id=C2-B8-D11-VD-049,muc=VD,loigiai={Sau $2$ giờ, hai máy bay ở $P_1=(-7;5;1)$ và $P_2=(8;-4;7)$ nên khoảng cách là $P_1P_2=\sqrt{342}=3\sqrt{38}$ (km).}]{Trong không gian $Oxyz$ (đơn vị: km), hai máy bay cùng xuất phát lúc $0$ giờ từ $A(-3;3;-1)$ và $B(4;0;3)$, bay thẳng với vận tốc không đổi lần lượt là $\vec{u}=(-2;1;1)$ và $\vec{v}=(2;-2;2)$ (km/h). Sau $2$ giờ, tính bình phương khoảng cách giữa hai máy bay (đơn vị: $\text{km}^2$).}{342}

\cauTLN[id=C2-B8-D11-VD-050,muc=VD,loigiai={Sau $2$ giờ, hai máy bay ở $P_1=(-1;-4;5)$ và $P_2=(0;3;7)$ nên khoảng cách là $P_1P_2=\sqrt{54}=3\sqrt{6}$ (km).}]{Trong không gian $Oxyz$ (đơn vị: km), hai máy bay cùng xuất phát lúc $0$ giờ từ $A(-1;-4;3)$ và $B(-2;3;3)$, bay thẳng với vận tốc không đổi lần lượt là $\vec{u}=(0;0;1)$ và $\vec{v}=(1;0;2)$ (km/h). Sau $2$ giờ, tính bình phương khoảng cách giữa hai máy bay (đơn vị: $\text{km}^2$).}{54}

\cauTLN[id=C2-B8-D11-VD-051,muc=VD,loigiai={Sau $1$ giờ, hai máy bay ở $P_1=(1;3;-3)$ và $P_2=(6;-5;-6)$ nên khoảng cách là $P_1P_2=\sqrt{98}=7\sqrt{2}$ (km).}]{Trong không gian $Oxyz$ (đơn vị: km), hai máy bay cùng xuất phát lúc $0$ giờ từ $A(-1;1;-2)$ và $B(4;-3;-4)$, bay thẳng với vận tốc không đổi lần lượt là $\vec{u}=(2;2;-1)$ và $\vec{v}=(2;-2;-2)$ (km/h). Sau $1$ giờ, tính bình phương khoảng cách giữa hai máy bay (đơn vị: $\text{km}^2$).}{98}

\cauTLN[id=C2-B8-D11-VD-052,muc=VD,loigiai={Sau $1$ giờ, hai máy bay ở $P_1=(5;-3;-4)$ và $P_2=(3;-1;-3)$ nên khoảng cách là $P_1P_2=\sqrt{9}=3$ (km).}]{Trong không gian $Oxyz$ (đơn vị: km), hai máy bay cùng xuất phát lúc $0$ giờ từ $A(4;-1;-4)$ và $B(4;0;-2)$, bay thẳng với vận tốc không đổi lần lượt là $\vec{u}=(1;-2;0)$ và $\vec{v}=(-1;-1;-1)$ (km/h). Sau $1$ giờ, tính bình phương khoảng cách giữa hai máy bay (đơn vị: $\text{km}^2$).}{9}

\cauTLN[id=C2-B8-D11-VD-053,muc=VD,loigiai={$AB=5$, $BC=7$ nên quãng đường bay là $5+7=12$ (km).}]{Trong không gian $Oxyz$ (đơn vị: km), một drone bay từ $A(-3;-1;0)$ đến $B(0;3;0)$ rồi bay thẳng đến $C(2;6;6)$. Tính tổng quãng đường drone đã bay (đơn vị: km).}{12}

\cauTLN[id=C2-B8-D11-VD-054,muc=VD,loigiai={$AB=3$, $BC=1$ nên quãng đường bay là $3+1=4$ (km).}]{Trong không gian $Oxyz$ (đơn vị: km), một drone bay từ $A(-1;-1;3)$ đến $B(0;1;5)$ rồi bay thẳng đến $C(0;1;6)$. Tính tổng quãng đường drone đã bay (đơn vị: km).}{4}

\cauTLN[id=C2-B8-D11-VD-055,muc=VD,loigiai={$AB=3$, $BC=11$ nên quãng đường bay là $3+11=14$ (km).}]{Trong không gian $Oxyz$ (đơn vị: km), một drone bay từ $A(-2;-2;2)$ đến $B(-1;-4;4)$ rồi bay thẳng đến $C(5;2;11)$. Tính tổng quãng đường drone đã bay (đơn vị: km).}{14}

\cauTLN[id=C2-B8-D11-VD-056,muc=VD,loigiai={$AB=11$, $BC=3$ nên quãng đường bay là $11+3=14$ (km).}]{Trong không gian $Oxyz$ (đơn vị: km), một drone bay từ $A(-2;-2;0)$ đến $B(4;4;7)$ rồi bay thẳng đến $C(2;6;8)$. Tính tổng quãng đường drone đã bay (đơn vị: km).}{14}

\cauTLN[id=C2-B8-D11-VD-057,muc=VD,loigiai={Vì $AM:MB=1:2$ nên $\overrightarrow{AM}=\dfrac{1}{3}\overrightarrow{AB}$, suy ra $M=(-1;-1;-3)$.}]{Trong không gian $Oxyz$ (đơn vị: km), hai trạm ở $A(0;-3;-5)$ và $B(-3;3;1)$. Một trạm chuyển tiếp $M$ nằm trên đoạn $AB$ sao cho $AM:MB=1:2$. Tính tổng ba tọa độ của trạm chuyển tiếp $M$.}{-5}

\cauTLN[id=C2-B8-D11-VD-058,muc=VD,loigiai={Vì $AM:MB=1:2$ nên $\overrightarrow{AM}=\dfrac{1}{3}\overrightarrow{AB}$, suy ra $M=(1;-1;1)$.}]{Trong không gian $Oxyz$ (đơn vị: km), hai trạm ở $A(2;0;4)$ và $B(-1;-3;-5)$. Một trạm chuyển tiếp $M$ nằm trên đoạn $AB$ sao cho $AM:MB=1:2$. Tính tổng ba tọa độ của trạm chuyển tiếp $M$.}{1}

\cauTLN[id=C2-B8-D11-VD-059,muc=VD,loigiai={Vì $AM:MB=2:1$ nên $\overrightarrow{AM}=\dfrac{2}{3}\overrightarrow{AB}$, suy ra $M=(1;-3;-3)$.}]{Trong không gian $Oxyz$ (đơn vị: km), hai trạm ở $A(3;-1;1)$ và $B(0;-4;-5)$. Một trạm chuyển tiếp $M$ nằm trên đoạn $AB$ sao cho $AM:MB=2:1$. Tính tổng ba tọa độ của trạm chuyển tiếp $M$.}{-5}

\cauTLN[id=C2-B8-D11-VD-060,muc=VD,loigiai={Vì $AM:MB=3:1$ nên $\overrightarrow{AM}=\dfrac{3}{4}\overrightarrow{AB}$, suy ra $M=(1;3;-4)$.}]{Trong không gian $Oxyz$ (đơn vị: km), hai trạm ở $A(-2;-3;-4)$ và $B(2;5;-4)$. Một trạm chuyển tiếp $M$ nằm trên đoạn $AB$ sao cho $AM:MB=3:1$. Tính tổng ba tọa độ của trạm chuyển tiếp $M$.}{0}
